\documentclass[11pt]{amsart}

\usepackage[margin=1.05in]{geometry}
\usepackage{amsmath,amssymb,amsthm}
\usepackage{booktabs}
\usepackage{array}
\usepackage{graphicx}
\usepackage[T1]{fontenc}
\usepackage{lmodern}
\usepackage{microtype}
\usepackage{xcolor}
\definecolor{linkblue}{rgb}{0.05,0.25,0.55}
\usepackage[colorlinks=true,linkcolor=linkblue,citecolor=linkblue,urlcolor=linkblue]{hyperref}

\newcommand{\F}{\mathbb{F}}
\newcommand{\Z}{\mathbb{Z}}
\newcommand{\Ot}{\widetilde{O}}
\newcommand{\Tt}{\widetilde{\Theta}}
\newcommand{\eps}{\varepsilon}
\newcommand{\Lp}[2]{L_{p}\!\left(#1,\,#2\right)}
\newcommand{\smoothpart}[2]{S_{#1}(#2)}
\newcommand{\osf}{\cite{OSF}}

\theoremstyle{plain}
\newtheorem{theorem}{Theorem}[section]
\newtheorem{lemma}[theorem]{Lemma}
\newtheorem{proposition}[theorem]{Proposition}
\theoremstyle{definition}
\newtheorem{definition}[theorem]{Definition}
\theoremstyle{remark}
\newtheorem{remark}[theorem]{Remark}

\begin{document}

\title[Elliptic curve primality proofs]{Elliptic curve primality proofs,\\ from very short to one-shot}
\author{Claude Code (Fable 5)}
\thanks{Written by Claude (Fable~5) under the direction of Andrew V.\ Sutherland, in the
\href{https://github.com/AndrewVSutherland2/OneShotFastECPP}{OneShotFastECPP} repository
(\texttt{reports/ecpp-varieties/}).  All verification-cost statements were checked against the reference
verifiers' code; all timing and count data are quoted from the repositories cited or measured on this
project's hardware.  Dickman $\rho$ values were computed by numerical solution of the delay differential
equation and cross-checked against published tables.}
\date{August 5, 2026}

\begin{abstract}
We survey algorithms and complexity bounds for four families of elliptic curve primality certificates ---
Pomerance proofs, one-shot ECPP (in its $n^4$-smooth and restricted $n^2$-smooth forms), short ECPP, and
classical ECPP --- as defined by the challenge repositories
\href{https://github.com/AndrewVSutherland/DANGER3}{DANGER3},
\href{https://github.com/AndrewVSutherland/OneShotPrimalityProofs}{OneShotPrimalityProofs}, and
\href{https://github.com/AndrewVSutherland/ShortPrimalityProofs}{ShortPrimalityProofs}.
For each family we give the problem statement, algorithms for finding certificates with heuristic
complexity analyses (treating the two order-supply engines, SEA point counting and the CM method, in
parallel throughout), the verification algorithm with a rigorous bit-complexity bound, and the best
results obtained to date.
\end{abstract}

\maketitle

\setcounter{tocdepth}{2}
\tableofcontents

\section{Introduction and summary}\label{sec:intro}

A \emph{primality certificate} for an integer $p$ is data whose validity (i) can be checked by an
efficient deterministic algorithm and (ii) implies, unconditionally, that $p$ is prime.
Pratt \cite{Pra75} observed that the classical Lucas criterion yields such a certificate for every prime
--- placing primality in NP --- but his certificates require the complete factorization of $p-1$
(recursively), so \emph{finding} them appears to require integer factorization, and verifying them takes
$O(\log^2 p)$ modular multiplications.  Two later developments changed the landscape.  Goldwasser and
Kilian \cite{GK86} showed how to use elliptic curves over $\F_p$, whose group orders vary over a Hasse
interval of width $4\sqrt p$, to \emph{construct} certificates in (heuristic, and for almost all primes
provable) polynomial time; and Pomerance \cite{Pom87} showed that elliptic curves make certificates
dramatically \emph{cheaper to verify}: every prime $p$ admits a proof requiring only
$(5/2+o(1))\log_2 p$ multiplications modulo $p$, i.e.\ $\Theta(n^2\log n)$ bit operations for an $n$-bit
prime --- a constant number of multiplications per bit, matching (up to the constant) Pepin's test for
Fermat primes and the Lucas--Lehmer test for Mersenne primes, but now valid for \emph{every} prime.

Pomerance's construction was an existence proof; he noted that finding such certificates ``may well take
exponential time.''  A cluster of 2026 challenge repositories of Sutherland
\cite{DANGER3,OneShot,Short} makes the gap between \emph{finding} and \emph{verifying} quantitative, by
fixing precise certificate formats (each with a reference verifier) and soliciting certificates for
ever-larger primes.  The four families, in the order treated here:

\begin{enumerate}
\item \textbf{Pomerance proofs} (repository DANGER3 \cite{DANGER3}): a single Montgomery curve with a
point of order $2^k$, where $2^k$ is the \emph{least} power of two exceeding the Goldwasser--Kilian
threshold ${\approx}\sqrt p$.  The certificate is three integers; verification is
${\approx}\,(5/2)\log_2 p$ multiplications, $O(n^2\log n)$ bit operations; the best known search is
$\Tt(\sqrt p)$ time.
\item \textbf{One-shot ECPP} (repository OneShotPrimalityProofs \cite{OneShot}): the power of two is
relaxed to any $n^4$-smooth order in a minimal window above the threshold
($n = \lceil\log_2 p\rceil$), with the certificate listing the prime factors in $(n^2, n^4)$.
Verification is $O(n^2(\log n)^2)$ bit operations; the best known searches are $p^{1/8+o(1)}$
(candidate orders, via SEA sampling) and $p^{1/4+o(1)}$ end-to-end (via the CM method).
\item \textbf{Restricted one-shot ECPP}: the sub-family with no listed primes --- the point's order must
be $n^2$-smooth outright.  The smoothness exponent doubles: $p^{1/4+o(1)}$ candidates (SEA),
$p^{1/2+o(1)}$ end-to-end (CM).  This interpolates between Pomerance proofs and one-shot ECPP;
certificates for $10^{20},\dots,10^{50}+\eps$ produced for this write-up appear in
\S\ref{sec:restricted}.
\item \textbf{Short ECPP} (repository ShortPrimalityProofs \cite{Short}): a \emph{chain} of such
certificates in which each level's point order is ($n^2$-smooth)${}\times{}$(one prime
$p_1 < \sqrt{p_0}$), and the next level certifies $p_1$.  Bit-lengths at least halve per level, so
certificates stay linear in $\log p$ and verification is $O(n^2\log n)$ bit operations ---
Pomerance-order (\S\ref{sec:shortverif}) --- while the free prime factor makes certificates
\emph{findable} in heuristic subexponential time $\Lp{1/3}{{\approx}1.3}$.
\item \textbf{Classical ECPP} (Goldwasser--Kilian, Atkin--Morain): the chain descends by an
$O(\log n)$-bit shrink per step with an essentially unconstrained cofactor.  Finding is heuristic
polynomial time --- $(\log p)^{4+\eps}$ for fastECPP --- at the price of quadratic-size certificates and
$O(n^3)$-bit-operation verification.
\end{enumerate}

The five formats form a hierarchy.  A Pomerance triple \emph{is} a restricted one-shot certificate
($m = 2^k$ is certainly $n^2$-smooth), which is a one-shot certificate with an empty prime list, and
also a zero-level short ECPP; a short ECPP is a maximally steep descent chain; classical ECPP relaxes
the descent rate and the order structure entirely.  Existence for \emph{every} prime $p > 3$ is
inherited down the whole hierarchy from Pomerance's theorem --- unconditionally, with no heuristics ---
while each relaxation strictly enlarges the set of valid certificates and thereby lowers the cost of
finding one.  Table~\ref{tab:summary} summarizes ($n = \lceil\log_2 p\rceil$; sizes exclude $p$ itself;
all verification bounds are rigorous and deterministic, all finding costs are heuristic; constants and
derivations are in the sections below).

\begin{table}[htbp]
\centering\scriptsize
\setlength{\tabcolsep}{4pt}
\resizebox{\textwidth}{!}{%
\begin{tabular}{@{}lllll@{}}
\toprule
certificate & size (bits) & verification (bit ops) & finding (heuristic) & exists for\\
\midrule
Pratt \cite{Pra75} & $O(n^2)$ & $O(n^3\log n)$ &
  $\Lp{1/3}{(64/9)^{1/3}{\approx}1.923}$ (NFS) & all $p$\\
AKS--Bernstein certificate \cite{AKS04,Berr05,Ber07,LP19} & $n^{1+o(1)}$ & $(\log p)^{4+o(1)}$ &
  $(\log p)^{2+o(1)}$ rand.;\ \ $(\log p)^{6}(2{+}\log\log p)^{c_0}$ det. & all $p$\\
Pomerance triple \cite{Pom87,DANGER3} & ${\approx}\,2n$ & $O(n^2\log n)$ &
  $\Tt(\sqrt p) = \Tt(2^{n/2})$ & all $p>3$\\
restricted one-shot (\S\ref{sec:restricted}) & ${\approx}\,2.5n$ & $O(n^2(\log n)^2)$ &
  $p^{1/4+o(1)}$ SEA;\ \ $p^{1/2+o(1)}$ CM & all $p>3$\\
one-shot, $n^4$-smooth \cite{OneShot} & ${\approx}\,2.5n$--$4n$ & $O(n^2(\log n)^2)$ &
  $p^{1/8+o(1)}$ SEA;\ \ $p^{1/4+o(1)}$ CM & all $p>3$\\
short ECPP \cite{Short} & $O(n)$, ${\approx}\,5n$--$7n$ & $O(n^2\log n)^{\dagger}$ &
  $\Lp{1/3}{c}$, $c \approx 1.31$ SEA; $1.39$ CM & all $p\ge 5$\\
ECPP \cite{GK86,AM93} & $\Theta(n^2/\log n)$ & $O(n^3)$ &
  $(\log p)^{6+o(1)}$ SEA;\ \ $(\log p)^{4+\eps}$ CM & all $p$ (heur.)\rlap{${}^{\ddagger}$}\\
\bottomrule
\end{tabular}}
\medskip
\caption{The certificate families.  The Pomerance-triple verification is exactly
$(\tfrac52+o(1))\,n$ multiplications $\bmod\ p$.
${}^{\dagger}$For fillers $m_i \le n^{O(1)}$ --- every published
certificate, and always arrangeable (\S\ref{sec:shortverif}); the format's worst case, already attained
by its $k=0$ stratum ($=$ the restricted one-shot), is $O(n^2(\log n)^2)$, and
Remark~\ref{rem:tightshort} tightens the format to guarantee $O(n^2\log n)$ worst-case.
${}^{\ddagger}$Heuristic for all $p$; provable for almost all $p$, and for all $p$ in expected
polynomial time (ZPP) via \cite{AH92}.  The AKS--Bernstein row is discussed in
Remark~\ref{rem:aks}; the deterministic finding entry is the Lenstra--Pomerance decision bound, whose
exponent $c_0$ is effectively computable but has never been made explicit.  The Pratt entry is the cost
of factoring $p-1$ (and recursively below, a geometrically decreasing cost) by the number field
sieve.}
\label{tab:summary}
\end{table}

\begin{remark}[the AKS family, in certificate form]\label{rem:aks}
Agrawal, Kayal, and Saxena \cite{AKS04} decide primality deterministically with no certificate at all,
in time $(\log p)^{21/2+o(1)}$; Lenstra and Pomerance \cite{LP19} reduce the exponent to $6$: their
Theorem~1 gives a deterministic algorithm running ``in time at most
$(\log n)^6\cdot(2+\log\log n)^{c_0}$'' for ``some effectively computable real number $c_0$'' ---
effectively computable in the sense that their proofs ``implicitly contain an algorithm for computing
$c_0$,'' which no one has carried out; so the sharpest citable deterministic bound is exactly
$(\log p)^6(2+\log\log p)^{c_0}$, and the $c$ in a formulation $O(n^6\log^{c+o(1)}n)$ is not known.
The \emph{certificate} form of the family descends from Berrizbeitia \cite{Berr05} through Cheng and
Mih\u{a}ilescu--Avanzi to Bernstein \cite{Ber07} (see \cite{Ber07} for the history), and is spelled
out in Appendix~\ref{app:aks}: every prime has a certificate with witness set $S = \{1\}$
(Theorem~\ref{thm:akscert}), of total size $(\log p)^{1+o(1)}$ bits --- i.e.\ $n^{1+o(1)}$,
essentially the same order as the elliptic certificates --- findable in \emph{expected time
$(\log p)^{2+o(1)}$} (the only randomness is sampling an irreducible $f$ and a witness $r$, both
dense) and verifiable in $(\log p)^{4+o(1)}$ bit operations, the cost of one $p^d$-th powering in a
$\Theta((\log p)^{3+o(1)})$-bit ring.  As a concrete instance,
$(1,\,57449,\,28724,\,16826,\,y,\,2,\,\{1\})$ is a certificate for $2^{1024}+643$.  Deterministically,
trying candidates for $f$ and $r$ in a fixed order heuristically finds a certificate in the same
$(\log p)^{2+o(1)}$ time (only the guarantee is lost); the proven deterministic route remains the
\cite{LP19} decision algorithm, whence the two finding entries in the table.  Note carefully that the
verification is quasi-\emph{quartic}, and structurally so (Appendix~\ref{app:aks}): this family does
\emph{not} subsume the quasi-quadratic verification of the elliptic formats --- see the closing open
problem of \S\ref{sec:synthesis}.
\end{remark}

Two engines supply candidate curve orders throughout, and they trade off differently in every section:
\textbf{SEA} point counting (sample a curve, compute its order; polynomial cost per candidate, orders
statistically ``random'') and the \textbf{CM method} (enumerate imaginary quadratic discriminants $D$
with $4p = t^2 + |D|v^2$; orders $p+1\pm t$ arrive at polylogarithmic amortized cost, but only
$\Tt(\sqrt B)$ of them from discriminants up to $B$, and committing to a winner costs a class polynomial
of degree $h(D) \approx \sqrt{|D|}$).  Where the searched event is polynomially rare, CM's cheap
candidates win; where it is exponentially rare, the $\sqrt B$ supply curve throttles CM and SEA's flat
polynomial cost per candidate wins asymptotically.  Measured crossovers appear in \S\ref{sec:oneshotfind}
and \S\ref{sec:restrictedfind}.

\section{Preliminaries}\label{sec:prelim}

\subsection{Cost model and notation}\label{sec:cost}

$p$ always denotes the integer whose primality is at issue (a ``probable prime'' during a search), and
$n = \lceil\log_2 p\rceil$ its bit length.  $M(t)$ denotes the bit cost of multiplying two $t$-bit
integers; $M(t) = O(t\log t)$ unconditionally \cite{HvdH21}.  A ``multiplication mod $p$'' costs
$O(M(n))$; following Pomerance we count squarings as multiplications and ignore additions.  Every
complexity bound below is stated in bit operations; multiplication counts appear only where an exact
constant is the natural statement (as in Pomerance's theorem), convertible at
$1 \text{ mult} = \Theta(M(n))$ bit operations.  $\Ot(\cdot)$ suppresses polylogarithmic factors.
The subexponential scale is
\[
  \Lp{\alpha}{c} \;=\; \exp\bigl( (c+o(1))\, (\ln p)^{\alpha} (\ln\ln p)^{1-\alpha} \bigr),
\]
so $\Lp{0}{c}$ is polynomial in $\log p$, $\Lp{1}{c}$ fully exponential, $\Lp{1/2}{\cdot}$ the
ECM/quadratic-sieve scale, and $\Lp{1/3}{\cdot}$ the number field sieve scale.  ``$p^{\theta+o(1)}$''
means $\exp((\theta+o(1))\ln p)$; the $o(1)$ terms here go to zero slowly (they hide $\ln\ln$ factors in
the exponent), which is why observed behavior at a few hundred bits is far friendlier than the exponent
suggests --- a point made quantitative in \S\ref{sec:oneshot} and \S\ref{sec:restricted}.

All \emph{verification} statements below are rigorous, deterministic, and unconditional: a verifier that
outputs True proves primality with no hypotheses.  All \emph{finding} statements are heuristic, resting
on the two randomness assumptions of \S\ref{sec:supply}--\ref{sec:smooth} (curve orders behave like
uniform Hasse-interval integers up to explicit local corrections; smooth-part probabilities follow the
Dickman distribution).  This asymmetry --- unconditional verification, heuristic search --- is intrinsic
to the subject and is the reason certificates matter at all: the expensive, unrigorous part need be run
only once, by anyone, while the cheap, rigorous part can be repeated by everyone.

\subsection{Curves over $\Z/p\Z$ and Montgomery arithmetic}\label{sec:mont}

For prime $p > 3$, an elliptic curve over $\F_p$ in Montgomery form is
\[
  E_{A,B}\colon\; B y^2 = x^3 + A x^2 + x, \qquad B(A^2-4) \ne 0,
\]
with the usual chord--tangent group law on its projective points; $\#E$ denotes the number of points
including the point at infinity $O$ \cite{Sil09,Mon87}.  Every Montgomery curve has the rational
2-torsion point $(0,0)$, so $4 \mid \#E$ always (the three $2$-torsion $x$-coordinates are $0$ and the
roots of $x^2+Ax+1$; at least one lies on $E$ or forces a point of order $4$).  Montgomery's $x$-only
arithmetic works with classes $(X : Z)$ representing $x = X/Z$ on the \emph{Kummer line}, identifying
$P$ with $-P$: doubling costs $5$ multiplications (with $c = (A+2)/4$ precomputed: $U = (X+Z)^2$,
$V = (X-Z)^2$, $X' = UV$, $Z' = (U-V)(V + c(U-V))$), differential addition costs $6$, and the Montgomery
ladder computes $(X:Z)$ of $[k]P$ in ${\approx}\,11\log_2 k$ multiplications.  Two features matter for
certificates.

\emph{Twist-agnosticism.}  The $x$-only formulas depend only on $A$, not $B$, and are simultaneously
valid on $E_{A,1}$ and its quadratic twist.  Every $x_0 \in \F_p$ is the $x$-coordinate of a point on
one of the two (which one depends on whether $x_0^3+Ax_0^2+x_0$ is a square).  All four certificate
formats therefore quantify over ``there exist $B, y_0$ such that $(x_0, y_0)$ lies on
$B y^2 = x^3+Ax^2+x$'': the verifier never tests whether $x_0$ is ``on the curve,'' because it is always
on one of the pair, and the order condition is what carries the proof.  This matches the set-up of
\cite{Pom87}, whose homogeneous model $y^2z = x^3+ax^2z+bxz^2$ is the general curve with a rational
point of order 2.

\emph{Soundness over $\Z/p\Z$ for composite $p$.}  The verifier runs the same formulas modulo a $p$ not
yet known to be prime.  The computation then commutes with reduction modulo every prime $\ell \mid p$
\emph{provided no ring division by a zero-divisor occurs}; verifiers enforce this with gcd guards
($\gcd(A^2-4, p) = 1$; a claimed point at infinity must be $(X : 0)$ with $\gcd(X, p) = 1$; a claimed
non-identity must have $\gcd(Z, p) = 1$).  A failed guard either rejects or exhibits a proper factor of
$p$.  This is the same discipline as Lenstra's ``pseudocurves'' in ECM \cite{Len87}, run in reverse: ECM
hopes for a zero-divisor, verification insists on their absence.

\subsection{The certification lemma}\label{sec:lemma}

All four families rest on one lemma, essentially Goldwasser--Kilian's \cite{GK86}, in the single-curve
form of \cite{Pom87}.  Write
\[
  L(p) \;=\; q + 1 + \lfloor 2\sqrt q\rfloor, \qquad q = \lfloor\sqrt p\rfloor,
\]
the integer form of $(p^{1/4}+1)^2$: by the Hasse bound \cite[V.1]{Sil09}, $L(p)$ is the largest
possible order of a point on any elliptic curve over any field $\F_\ell$ with $\ell \le \sqrt p$.

\begin{lemma}[elliptic certification]\label{lem:cert}
Let $p > 3$ be an odd integer, and suppose that for every prime $\ell$ dividing $p$, the reduction
modulo $\ell$ of some fixed pair $(A, x_0)$ yields a point of order exactly $m$ on a Montgomery curve
$By^2 = x^3+Ax^2+x$ over $\F_\ell$, for one and the same integer $m > L(p)$.  Then $p$ is prime.
\end{lemma}

\begin{proof}
If $p$ were composite it would have a prime divisor $\ell \le \sqrt p$.  The point's order $m$ divides
$\#E(\F_\ell) \le \ell + 1 + 2\sqrt\ell \le L(p) < m$, a contradiction.
\end{proof}

Establishing ``order exactly $m$ modulo every prime divisor'' from $x$-only data is the verifier's whole
job: it checks that $[m](x_0{:}1)$ is the genuine point at infinity ($Z \equiv 0$ and
$\gcd(X, p) = 1$, which forces $[m]P = O$ modulo \emph{every} $\ell \mid p$), and that
$[m/q](x_0{:}1)$ is a unit-$Z$ point (hence ${}\ne O$ modulo every $\ell \mid p$) for each maximal
proper divisor $m/q$ of $m$, $q$ running over the distinct primes of $m$.  This requires \emph{knowing
the prime factorization of $m$} --- which is why every format constrains $m$'s arithmetic shape, and why
the differences between the formats are exactly differences in how $m$ may be built.

\begin{remark}[why ``modulo every prime divisor'' and not ``order in $E(\Z/p\Z)$'']\label{rem:crt}
The order of a point over the ring $\Z/p\Z$ is, by CRT, only the \emph{least common multiple} of its
orders modulo the primes dividing $p$; for composite $p$ the required order can be split between factors
with every local order strictly smaller than $m$.  A concrete counterexample from \cite{Short}:
$p_0 = 4410667997551 = 2098153\cdot 2102167$ with the tuple
$(p_0,\, 1365834658413,\, 107710304518,\, 4200232,\,\dots)$ has a point of order exactly
$8\cdot 525029$ in $E(\Z/p_0\Z)$ --- order $8$ modulo the first prime and $525029$ modulo the second ---
and would pass a verifier that read ``order'' in $\Z/p_0\Z$.  The gcd guards are what rule this out:
they force the ladder's intermediate values to behave identically modulo every prime divisor, so the
order is $m$ locally everywhere, and the reference verifiers reject the example (exposing both factors
of $p_0$ along the way).
\end{remark}

\subsection{The supply of curve orders}\label{sec:supply}

Fix a prime $p$.  Which integers occur as $\#E$ for $E/\F_p$, and how are they distributed?  Three
classical facts calibrate every heuristic in this write-up.

\begin{itemize}
\item \textbf{Hasse} \cite{Sil09}.  $\#E = p + 1 - t$ with $|t| \le 2\sqrt p$: orders live in the
\emph{Hasse interval} of length $4\sqrt p$ centered at $p+1$.
\item \textbf{Deuring} \cite{Deu41,Wat69}.  Conversely \emph{every} such $t$ occurs for some $E/\F_p$
($p$ prime).  Quantitatively, the number of curves with trace $t$, counted up to isomorphism and
weighted naturally, is the Hurwitz class number $H(4p-t^2)$ \cite{Deu41,Scho87}, of typical size
$\sqrt p$ up to logarithmic factors: a random curve has an essentially uniformly distributed trace,
modulated by a smooth semicircular-type density and by class number fluctuations.
\item \textbf{Group structure.}  $E(\F_p) \cong \Z/n_1 \times \Z/n_2$ with $n_1 \mid n_2$ and
$n_1 \mid p-1$; a point of order $m$ exists iff $m \mid n_2$ (the exponent).  For a curve chosen at
random, $n_1$ is $1$ or very small with high probability, and one can always move within the isogeny
class to make any given Sylow subgroup cyclic (\S\ref{sec:cmengine}).
\end{itemize}

\emph{The randomness heuristic.}  For searching we model the order of a random curve as a uniform
random integer in the Hasse interval, corrected by known local biases: $\#E$ is even with probability
${\approx}\,2/3$ over all curves, and exactly $1$ for Montgomery curves (which force $4 \mid \#E$); more
generally $\Pr[\ell \mid \#E] \approx \ell/(\ell^2-1)$ for small primes $\ell$
\cite[Prop.~1.14]{Len87}.  These biases \emph{help} every smoothness-hunting search below by a constant
factor, and are one ingredient of the $O(1)$ ``fudge factors'' between the clean Dickman model and
observed candidate counts (measured in \S\ref{sec:oneshotfind} and \S\ref{sec:restrictedresults} to lie
between about $1/2$ and $20$).

\subsection{Smooth numbers}\label{sec:smooth}

An integer is \emph{$y$-smooth} if all its prime factors are ${}\le y$.  With
$\Psi(x, y) = \#\{\,m \le x : m \text{ is } y\text{-smooth}\,\}$, the fundamental estimate
\cite{Dic30,CEP83,Hil86} is
\[
  \Psi(x, y) \;=\; x\,\rho(u)^{1+o(1)}, \qquad u = \ln x / \ln y,
\]
uniformly for $u < y^{1-\eps}$, where the Dickman function $\rho$ satisfies $\rho(u) = 1$ on $[0,1]$ and
$u\rho'(u) = -\rho(u-1)$; asymptotically
$\rho(u) = u^{-u(1+o(1))} = \exp(-u(\ln u + \ln\ln u - 1 + o(1)))$.  See \cite{Gran08} for a survey.
Two consequences are used constantly.

\emph{Fully smooth orders.}  $\Pr[\text{a random Hasse-interval integer is } y\text{-smooth}] \approx
\rho(u)$ with $u = \ln p/\ln y$.  With $y$ polylogarithmic in $p$,
$u \approx \ln p/(c\ln\ln p)$ and $\rho(u) = p^{-1/c+o(1)}$: \emph{polylog-smooth orders are a positive
power of $p$ rare}.  This single computation drives all the exponents $1/2$, $1/4$, $1/8$ below.

\emph{Smooth parts.}  Write $\smoothpart{y}{N}$ for the largest $y$-smooth divisor of $N$.  The
certificates of \S\S\ref{sec:oneshot}--\ref{sec:short} need only
$\smoothpart{y}{\#E} > L(p) \approx \sqrt p$ --- half the bits --- not full smoothness.  For
$N \approx p$ random, $\Pr[\smoothpart{y}{N} > \sqrt N] = \rho(u/2)^{1+o(1)}$: up to factors absorbed in
the $o(1)$ (an $e^{-\gamma}$ from the rough cofactor, an integral over how far above $\sqrt N$ the
smooth part lands), demanding a smooth \emph{part} above $\sqrt N$ costs what full smoothness would cost
for $\sqrt N$.  Halving the bits that must be smooth halves $u$, and since $\rho$ is doubly exponential
in $u$ this is an enormous relaxation: at 256 bits with $y = n^4$ it is the difference between
$\rho(8) \approx 3\times 10^{-8}$ and $\rho(4) \approx 5\times 10^{-3}$.

\subsection{Two engines for curve orders: SEA and CM}\label{sec:engines}

Every finder below is a filter over candidate pairs (curve, order).  There are exactly two known ways to
manufacture such pairs, and their cost structures are dual.

\subsubsection{Point counting: Schoof--Elkies--Atkin}\label{sec:seaengine}
Schoof's algorithm \cite{Scho85} computes $\#E$ for arbitrary $E/\F_p$ in deterministic polynomial time
by computing the trace $t \bmod \ell$ for the primes $\ell = O(\log p)$, working in the $\ell$-division
polynomials of degree $\Theta(\ell^2)$; with fast arithmetic it runs in $\Ot(n^5)$.  The Elkies--Atkin
improvements \cite{Scho95,Elk98} replace division polynomials, for the ${\approx}\,$half of primes
$\ell$ that are ``Elkies primes,'' by degree-$\Theta(\ell)$ kernel polynomials of rational
$\ell$-isogenies located via the modular polynomial $\Phi_\ell(X, j(E))$, giving the SEA algorithm with
heuristic cost $\Ot(n^4)$ (the heuristic being the supply of Elkies primes; see \cite{SS14} for what is
provable on average).  As an order \emph{supply}, SEA has two decisive properties: the cost per
candidate is a \emph{flat polynomial} --- independent of how many candidates one consumes --- and each
quadratic twist comes free ($\#E + \#E^{\mathrm{tw}} = 2p+2$), so one count yields two candidate orders.
Its constant is large in practice: a stock PARI/GP \texttt{ellcard} costs 7--14 core-seconds per curve
at 384--448 bits \osf.

\subsubsection{Complex multiplication: prescribed traces from small discriminants}\label{sec:cmengine}
The CM method inverts the problem: instead of choosing a curve and learning its order, choose the order
and construct a curve.  If $D < 0$ is an imaginary quadratic discriminant with
\[
  4p = t^2 + |D|v^2
\]
solvable in integers (equivalently, $p$ splits suitably in the ring class field of discriminant $D$),
then there are curves $E/\F_p$ with CM by the order of discriminant $D$ and trace $\pm t$, hence orders
$p+1-t$ and $p+1+t$; for $D = -3, -4$ the extra units give six resp.\ four orders \cite{AM93,Cox13}.
The machinery has three separately priced stages.

\begin{enumerate}
\item \emph{Which $D$ are usable, and what are $t, v$?}  Testing one $D$ costs a Legendre symbol, a
square root of $D \bmod p$, and Cornacchia's algorithm \cite{Cor08}, \cite[\S2.3.4]{CP05} ---
polylogarithmic.  Batched over many $D$, the square roots reduce to one square root per prime factor of
the discriminants plus one modular multiplication per $D$ (Shallit's trick; credited in
\cite{Mor07,FKMW04}), and restricting to discriminants all of whose prime discriminant factors are
residues on $p$ boosts the solvability probability by genus theory
($\Pr = 2^{\#\text{factors}-1}/h(D)$ for such $D$).  The engine of \osf{} scans 5--20 million
discriminants per second at 128--384 bits this way.
\item \emph{How many orders does a budget of discriminants buy?}  This is the binding constraint.  The
solvability probability for a single $D$ is ${\approx}\,c/h(D) \approx c'/\sqrt{|D|}$ (up to
$L(1, \chi_D)$ factors), so discriminants up to $B$ yield only
\[
  \#\{\text{usable } (D, t)\} \;\approx\; \sum_{|D|\le B} \Theta(1/\sqrt{|D|}) \;\approx\; \Tt(\sqrt B)
\]
candidate orders, at scan cost $\Tt(B)$.  Empirically the constant is ${\approx}\,0.5$: the run of
\osf{} on $10^{110}+7$ measured a pool of 39{,}878 orders at
$B = 8\times 10^9$, i.e.\ ${\approx}\,0.45\sqrt B$.  \emph{Producing $K$ candidate orders by CM costs
$\Tt(K^2)$.}
\item \emph{Committing to a winner.}  Only the finally selected $D$ requires the expensive step: the
Hilbert class polynomial $H_D$ (or a smaller class invariant) of degree $h(D) \approx \sqrt{|D|}$,
computable in quasi-linear time $\Ot(|D|)$ by the CRT method \cite{Suth11,Enge09,Suth12}, one root of it
$\bmod\ p$, the $j$-invariant, the right twist, and a point of the desired order.  Isogeny-volcano
navigation \cite{Kohel96,FM02,Suth13} fixes group-structure obstructions (\S\ref{sec:oneshotfind}).
\end{enumerate}

The duality, which recurs in every section: \textbf{SEA pays $\mathrm{poly}(n)$ per candidate forever;
CM pays polylog per candidate but quadratically in the number of candidates.}  If a search needs $K$
candidates, SEA costs $\Ot(K)\,\mathrm{poly}(n)$ and CM costs $\Ot(K^2) + \Ot(K\,\mathrm{polylog})$.
CM wins when $K$ is modest (all of classical ECPP, one-shot to ${\approx}\,450$ bits); SEA wins when $K$
explodes (one-shot beyond ${\approx}\,500$ bits, the restricted variant, and --- in the degenerate
single-candidate sense of \S\ref{sec:pomerance} --- never for Pomerance triples, where the admissible
orders are known in advance and neither engine is needed).

\subsection{The prover's factoring toolbox}\label{sec:toolbox}

Several finders must partially factor candidate orders (never as part of any certificate's validity ---
only to \emph{recognize} a usable order).  Trial division to $y$ costs $\pi(y)$ divisions per candidate
naively, and amortizes to polylogarithmic per candidate over large batches via product/remainder trees
\cite{Bern04}; ECM \cite{Len87} finds a prime factor $q \le y$ of $N$ in heuristic expected time
$L_y(1/2, \sqrt 2)\cdot\Ot(M(\log N))$, so ``peel everything ${}\le y$ from $N$'' costs
$L_y(1/2, \sqrt 2)$ up to polynomial factors; the number field sieve \cite{LL93} factors $N$ completely
in heuristic $L_N(1/3, (64/9)^{1/3} \approx 1.923)$.  The interplay ``how much peeling per candidate
vs.\ how rare is a peelable candidate'' is what produces the $L(1/3)$ finding costs of
\S\ref{sec:short}.

\section{Pomerance proofs}\label{sec:pomerance}

\subsection{The problem}

Pomerance's 1987 theorem \cite[Thm.~1]{Pom87}: \emph{for every prime $p$ there is a proof of its
primality whose verification requires $(5/2+o(1))\log_2 p$ multiplications $\bmod\ p$.}  His certificate
exhibits a point of order $2^k$ on an elliptic curve, with $2^k$ just above the threshold $L(p)$ of
\S\ref{sec:lemma}, so that verification is nothing but $k \approx \frac12\log_2 p$ successive doublings.
The DANGER3 repository \cite{DANGER3} (the data challenge of the April 2026 BIRS workshop \emph{DANGER:
Data, Numbers, and Geometry}) fixes the modern Montgomery-form refinement:

\begin{definition}[{\cite{DANGER3}}]\label{def:pomtriple}
A \emph{Pomerance triple} is a triple of integers $(p, A, x_0)$ with $p$ a positive odd integer,
$0 \le A, x_0 < p$, $A \ne \pm 2 \bmod p$, such that there exist $B$ and $y_0$ making $(x_0, y_0)$ a
point of order $2^k$ on the Montgomery curve $By^2 = x^3 + Ax^2 + x$ modulo every prime divisor of $p$,
where $k$ is the \emph{least} integer with $2^k > q + 1 + 2\sqrt q$, $q = \lfloor\sqrt p\rfloor$.
Operationally: applying the Montgomery doubling law $k-1$ times to $(x_0 : 1)$ gives a point with
$Z$-coordinate coprime to $p$, and the $k$-th doubling gives $Z \equiv 0 \pmod p$.
\end{definition}

By Lemma~\ref{lem:cert} a valid triple proves $p$ prime ($2^k > L(p)$), and by
\cite[Lemmas 2.1--2.2]{Pom87} the doubling-chain behavior ($Z \ne 0$ at step $k-1$, $Z = 0$ at step $k$)
characterizes ``order exactly $2^k$'' on the Kummer line.  The certificate is \emph{two} integers beyond
$p$ itself --- about $2n$ bits --- the shortest of the elliptic-curve formats.

\subsection{Existence: every prime has one}\label{sec:pomexist}

This is the part that uses the full strength of Deuring's theorem, and it is worth sketching because
every other format inherits its existence statement from this argument \cite[Thm.~3.3]{Pom87}.

\begin{enumerate}
\item Minimality puts $2^k$ in $(L(p), 2L(p)]$, which is below the Hasse interval's width $4\sqrt p$
once $p > 34$, so the interval $(p+1-2\sqrt p,\, p+1+2\sqrt p)$ contains a multiple $m = 2^k c$ of
$2^k$.
\item By Deuring \cite{Deu41}, some $E/\F_p$ has $\#E = m$ ($m$ is even, so $E$ can be taken with a
rational $2$-torsion point, i.e.\ in the model of \S\ref{sec:mont}).
\item $E(\F_p)$ is a product of at most two cyclic groups.  If some curve of order $m$ has a point of
order $2^k$, done (Pomerance's ``type 1'' certificate).  Otherwise the $2$-Sylow subgroup splits;
Pomerance's ``type 2'' certificate exhibits two points of orders $2^{k_1}, 2^{k_2}$ with
$k_1+k_2 = k$ generating a subgroup of order $2^k$, which suffices for the same contradiction.  For the
single-point Montgomery format one instead moves within the isogeny class: descending the $2$-isogeny
volcano to its floor \cite{Kohel96,Suth13} makes the $2$-Sylow subgroup cyclic, so a point of order
$2^k$ exists there, and the floor curve is Montgomery-representable.  This is the ``minor refinement''
recorded in \cite{DANGER3}: \emph{Pomerance triples exist for all primes $p > 3$}.
\end{enumerate}

Empirically the repository verifies this exhaustively far beyond the theorem's needs: all 43 million
triples with $p < 2^{14}$, and at least one triple for \emph{every} prime $3 < p < 2^{32}$
(203{,}280{,}219 of them).

\subsection{Verification, exactly}\label{sec:pomverif}

The reference verifier \texttt{vpp.py} \cite{DANGER3} computes $k = {}$bit-length of
$q+1+\lfloor 2\sqrt q\rfloor$, checks $\gcd(A^2-4, p) = 1$, and runs $k$ Montgomery doublings from
$(x_0 : 1)$ at 5 multiplications each (\S\ref{sec:mont}), accepting iff $Z_k \equiv 0$ and
$\gcd(Z_{k-1}, p) = 1$.  Total:
\[
  5k + O(1) \;=\; (5/2 + o(1))\log_2 p \ \text{ multiplications mod } p,
\]
i.e.\ $O(n\,M(n)) = O(n^2\log n)$ bit operations, plus two gcds.  This realizes Pomerance's constant: his Theorem~3.3 proves
$(7/2)\log_2 p + O(1)$ multiplications with the 7-multiplication doubling of his general two-parameter
model, and the remarks following it reduce $7/2$ to $5/2$ --- precisely the 5-multiplication Montgomery
doubling (Montgomery's paper \cite{Mon87} appears in the same volume of \emph{Mathematics of
Computation}).  For calibration: a \emph{single} strong-pseudoprime test --- which proves nothing ---
costs $(1+o(1))\log_2 p$ multiplications; a Pomerance triple proves primality outright for $2.5\times$
that.  The repository also provides \texttt{lean\_vpp.py}, which emits a Lean~4 script certifying a
given triple, so verification can even be delegated to a proof assistant's kernel.

\subsection{Finding certificates: $\Tt(\sqrt p)$, and why}\label{sec:pomfind}

Here the supply engines of \S\ref{sec:engines} are both nearly useless, which is what makes this family
the hardest to find and the sharpest illustration of the find/verify gap.

\emph{The counting heuristic.}  A curve is usable iff $\#E$ is one of the multiples of $2^k$ in the
Hasse interval --- two to four integers, \emph{computable in advance from $p$ alone} since
$2^k \in (L(p), 2L(p)] \approx (1\dots 2)\sqrt p$ --- and iff its group has a cyclic-enough $2$-Sylow
subgroup with the sampled point generating its full 2-part.  Under \S\ref{sec:supply} each admissible
order is hit with probability $\Tt(1/\sqrt p)$, and the structure conditions cost only a constant, so
the per-curve success probability is $\Tt(1/\sqrt p)$, the hidden constant wobbling by a factor
${\approx}\,4$ with the position of $\sqrt p$ between consecutive powers of two.  So $\Tt(\sqrt p)$
curves must be sampled, and the algorithmic content is making the per-candidate test essentially free.

\emph{The 2-Sylow projection test} (Ruehle's implementation \cite{Rue26}, used for the $10^{19}$ and
$10^{21}$ records).  Let $N_1,\dots,N_j$ ($j \le 4$) be the admissible orders and
$c_i = N_i/2^{v_2(N_i)}$ their odd parts.  For a random pair $(A, x_0)$: one Montgomery ladder computes
$[\prod c_i](x_0 : 1)$, which lands in the 2-Sylow subgroup \emph{whenever the curve's order is
admissible}; then repeated doublings with early abort either die (reject) or reach $Z = 0$ at exactly
depth $k$ (accept, and the affine $x$-coordinate of the point $k$ steps back is the certificate's
$x_0$).  Without the projection, a random point has 2-power order with probability only $\Tt(1/p)$ ---
the projection is what turns $p$ trials into $\sqrt p$ trials.  Per-candidate cost: one ladder of $O(n)$
bits plus $O(k)$ cheap doublings, i.e.\ $O(n)$ multiplications on numbers that fit in one or two machine
words at challenge sizes --- a few hundred nanoseconds on a CPU core, a few nanoseconds amortized on a
GPU.  Sampling curve families with prescribed 2-power torsion (via the modular curve $X_1(16)$)
conditions $v_2(\#E)$ and buys a further constant factor, as observed in the $10^{23}$ search
\cite{McL26} --- explicitly a constant-factor gain, not an asymptotic one.

\emph{Why neither engine helps.}  SEA point counting per candidate ($\Ot(n^4)$ each) is strictly worse
than the $\Ot(n)$-multiplication direct test, and no partial-information shortcut is useful:
$t \bmod 2^j$ for $j$ up to $k \approx n/2$ is exactly as hard as the event being tested.  The CM route
would need a small discriminant $D$ with $4p = t^2+|D|v^2$ and $t \equiv p+1 \pmod{2^k}$; only $O(1)$
values of $t$ satisfy the congruence inside the Hasse range, and for fixed $t$ the probability that
$4p-t^2$ has squarefull part large enough to leave $|D| \le X$ is ${\approx}\,\sqrt{X/p}$ --- total
expected work $p^{1/2-o(1)}$ regardless.  Isogeny walks preserve the order, torsion families fix only
$O(1)$ bits of $v_2(\#E)$, and the admissible set has no multiplicative structure to sieve on.  Nothing
better than $\Tt(\sqrt p)$ is known, and the DANGER3 progression (below) is the empirical record of that
$\sqrt p$ wall being pushed one decimal digit --- half a bit of $\sqrt p$ --- at a time.

\subsection{Results to date}\label{sec:pomresults}

\cite{DANGER3} lists example triples for $p = 10^{12}+39$ through $10^{18}+3$, then a challenge log in
which every entry is a human--LLM collaboration (Table~\ref{tab:danger}).  Candidates are $(A, x_0)$
pairs tested by the projection method; $\sqrt p$ is shown for comparison with the $\Tt(\sqrt p)$
heuristic.

\begin{table}[htbp]
\centering\footnotesize
\begin{tabular}{@{}lllll@{}}
\toprule
$p$ & $\sqrt p$ & candidates & time / hardware & found by (date)\\
\midrule
$10^{19}+51$  & $3.2\times 10^{9}$  & $3.7\times 10^{9}$  & 400 s, multithreaded C & F.~Ruehle + Claude Opus 4.6 (Apr 2026)\\
$10^{20}+39$  & $1.0\times 10^{10}$ & $1.2\times 10^{9}$  & 7 h                    & V.~Jejjala + GPT 5.4 Pro (Apr 2026)\\
$10^{21}+117$ & $3.2\times 10^{10}$ & $5.3\times 10^{10}$ & 16 h                   & F.~Ruehle + Claude Opus 4.6 (Apr 2026)\\
$10^{22}+9$   & $1.0\times 10^{11}$ & $5.9\times 10^{10}$ & 16 h                   & A.~McLain + GPT 5.5 Codex (May 2026)\\
$10^{23}+117$ & $3.2\times 10^{11}$ & $3.1\times 10^{10}$ & ---                    & A.~McLain + GPT 5.5 Codex (Jun 2026)\\
$10^{24}+7$   & $1.0\times 10^{12}$ & $1.8\times 10^{10}$ & 6 h, Apple M3 laptop   & J.~Shi + Claude Fable 5 (Jun 2026)\\
$10^{25}+13$  & $3.2\times 10^{12}$ & $2.0\times 10^{11}$ & ---                    & A.~McLain + GPT 5.5 Codex (Jun 2026)\\
$10^{26}+67$  & $1.0\times 10^{13}$ & $1.4\times 10^{11}$ & 45 min, RTX 6000 GPU   & A.~McLain + GPT 5.5 Codex (Jun 2026)\\
\bottomrule
\end{tabular}
\medskip
\caption{The DANGER3 challenge log for Pomerance triples.}
\label{tab:danger}
\end{table}

The observed candidate counts sit between $0.01\sqrt p$ and $1.7\sqrt p$ --- the spread reflecting the
$2^k/\sqrt p$ wobble, torsion conditioning, twist usage, and single-draw geometric variance --- and the
per-candidate cost has been driven from CPU-microseconds to GPU-nanoseconds.  The open challenge as of
this writing is $p = 10^{27}+103$ ($\sqrt p \approx 3.2\times 10^{13}$ candidates --- roughly ten
RTX-6000-hours at the $10^{26}$ rate: within reach).  Note the striking asymmetry these records
quantify: a certificate that takes GPU-\emph{hours} to find is verified in ${\approx}\,220$
multiplications of 90-bit numbers --- microseconds.

\section{One-shot ECPP}\label{sec:oneshot}

\subsection{The problem}

The OneShotPrimalityProofs repository \cite{OneShot} (part of the DARPA expMath program) generalizes the
Pomerance triple by relaxing ``order a power of two'' to ``order $n^4$-smooth,'' keeping the single
curve --- one shot, no recursion --- and the quasi-quadratic verification.

\begin{definition}[{\cite{OneShot}}]\label{def:oneshot}
A \emph{one-shot ECPP} is a tuple $(p, A, x_0, m, q_1, \dots, q_k)$ with $p > 1$ odd,
$0 \le A, x_0 < p$, $A \ne \pm 2 \bmod p$, such that $(x_0, y_0)$ is a point of order $m$ on
$By^2 = x^3+Ax^2+x$ (for some $B, y_0$), where, with $n = \lceil\log_2 p\rceil$ and $L = L(p)$:
$m$ is \emph{$n^4$-smooth}; $L < m < L\cdot r$ where $r$ is the least prime divisor of $m$; and
$q_1 < \dots < q_k$ are the prime divisors of $m$ in $(n^2, n^4)$.
\end{definition}

Every Pomerance triple is the case $k = 0$, $m = 2^k$ minimal --- so one-shot certificates exist for
every prime $p > 3$, unconditionally.  The minimality window $m < L\cdot r$ generalizes ``least power of
two above $L$'': it forces $m/r \le L$, i.e.\ removing the smallest prime of $m$ would fall out of the
window, keeping $\log m = \frac12\log p + O(\log n)$ and the certificate canonical-small.  Since each
listed prime exceeds $n^2$, $k < n/(4\log_2 n) + O(1)$, so the whole certificate is
${\approx}\,2.5n$--$4n$ bits.

\subsection{Verification}\label{sec:oneshotverif}

The reference verifier \texttt{voneshot.py} \cite{OneShot} establishes ``order exactly $m$ modulo every
prime divisor of $p$'' and the window conditions, as follows.

\begin{enumerate}
\item \emph{Window:} $L < m \le p+1+\lfloor 2\sqrt p\rfloor$ (a point's order cannot exceed the Hasse
bound if $p$ is prime; soundness rests on Lemma~\ref{lem:cert}, not on this check alone).
\item \emph{Nonsingularity:} $\gcd(A^2-4, p) = 1$.
\item \emph{Factor $m$:} sieve the primes to $n^2$ and extract $m$'s prime divisors below $n^2$ by
batched remainder trees \cite{Bern04}; strip them to get the $n^2$-rough part $m'$; then check that the
listed $q_i$ are strictly increasing, lie in $(n^2, n^4)$, and \emph{exactly exhaust $m'$}.  Two design
points deserve emphasis:
\begin{itemize}
\item \emph{The $q_i$ need no primality tests.}  Any divisor of $m'$ below $n^4$ is automatically prime:
a composite one would have two factors ${}> n^2$, exceeding $n^4$.  The certificate format chose the
exponent $4$ precisely to make the listed primes self-certifying.
\item \emph{The $q_i$ are validated against $m'$, not $m$.}  Validating against $m$ would admit a
``fold attack'': a repeated small prime $c \le n^2$ of $m$ could be folded into a fake ``large prime''
$cq$ dividing $m$ but not $m'$, corrupting the exact-order argument.  (A concrete attack tuple rejected
by this check is in the verifier's test suite.)
\end{itemize}
\item \emph{Minimality:} $m < L\cdot r$ with $r$ the least prime found.
\item \emph{$[m]P = O$ genuinely:} the ladder value $(X : Z)$ at $m$ must have $Z \equiv 0$ \emph{and}
$\gcd(X, p) = 1$ --- the second condition rejects the degenerate $(0 : 0)$ that arises if the ladder ran
past the local order modulo some prime divisor of $p$, and is what makes the step sound over $\Z/p\Z$
(\S\ref{sec:lemma}).
\item \emph{Exact order:} with $R = \prod q$ over the distinct primes $q \mid m$, compute
$Q = [m/R]P$ once, then verify $[m/q]P \ne O$ for every $q$ by a balanced divide-and-conquer over the
prime set: split the primes in halves, multiply $Q$ by each half's product, recurse on the opposite
half.  Each of the $O(\log n)$ recursion levels costs $O(n)$ group operations in total.
\end{enumerate}

\begin{proposition}[verification cost, rigorous]\label{prop:oneshotverif}
$O(n\log n)$ group operations, each $O(1)$ multiplications $\bmod\ p$, plus $\Ot(n^2)$ bit operations of
sieving and trial division: $O(n^2(\log n)^2)$ bit operations and $O(n^2)$ bits of memory ---
quasi-quadratic, a factor ${\approx}\,\log\log p$ above Pomerance triples (the depth of the product
tree), against quasi-cubic for a classical ECPP chain (\S\ref{sec:ecpp}).  Measured \osf: 8--14 ms in
pure Python at 256--333 bits.  This is the worst case over certificates; \S\ref{sec:shortverif} gives
the accounting parametrized by the number of distinct primes of $m$, which recovers $O(n^2\log n)$ when
that number is bounded.
\end{proposition}

\subsection{Finding certificates}\label{sec:oneshotfind}

Since $m$ divides $\#E$ and every divisor of the $n^4$-smooth part of $\#E$ is reachable as a point
order (up to group-structure corrections handled below), a curve is usable precisely when
\[
  \smoothpart{n^4}{\#E} \;>\; L \;\approx\; \sqrt p,
\]
in which case the minimal in-window divisor $m$ of the smooth part is extracted greedily (strip the
largest primes while staying above $L$; the greedy endpoint automatically satisfies $m < L\cdot r$).  By
\S\ref{sec:smooth} the per-candidate probability is $\rho(u/2)^{1+o(1)}$ with
$u = \ln p/\ln n^4 = \ln p/(4\ln\ln p)(1+o(1))$, and
\[
  \rho(u/2) = \exp\bigl(-(u/2)(\ln u)(1+o(1))\bigr) = p^{-1/8+o(1)},
\]
so $p^{1/8+o(1)}$ \emph{candidate orders} are needed.  The $o(1)$ is strongly negative at practical
sizes: $1/\rho(u/2)$ is only ${\approx}\,200$ at 256 bits, $1.6\times 10^4$ at 384, $1.6\times 10^5$ at
448, ${\approx}\,3\times 10^8$ at 640 bits.  What the two engines pay per candidate then separates
them.

\subsubsection{SEA sampling: $p^{1/8+o(1)}$ end-to-end}
The repository's \texttt{oneshot.gp} and the race harness \texttt{sea\_supply.py} \osf{} sample random
Montgomery coefficients $A$, compute $\#E$ by SEA (both twists per count), batch-test the smoothness
gate, and assemble the winner.  Per candidate: half an SEA count ($\Ot(n^4)$ heuristic) plus
polylogarithmic gate work, so the total is $p^{1/8+o(1)}\cdot\mathrm{poly}(n) = p^{1/8+o(1)}$.  The
measured per-candidate cost with stock PARI (\texttt{ellcard}) is flat: 7.2 / 10.1 / 13.7 core-seconds
at 384 / 416 / 448 bits.

\subsubsection{The CM engine: $p^{1/4+o(1)}$ end-to-end, far cheaper per candidate today}
\label{sec:cmoneshot}
The OneShotFastECPP engine \osf{} is a fastECPP-style pipeline specialized to the one-shot problem ---
unusual in that its $p$ is small (hundreds of bits) while its discriminants are huge ($|D|$ up to
$10^{11}$), the opposite of classical fastECPP, which relocates all the costs:
\begin{enumerate}
\item \emph{Discriminant scan}: enumerate every fundamental $D$ with $|D| < B$ all of whose prime
discriminant factors are residues on $p$ (genus theory maximizes the solvability rate and keeps
$\sqrt D \bmod p$ computable multiplicatively from a factor base of $\sqrt{q^*}$ values --- one
Tonelli--Shanks per factor-base prime, one modular multiplication per $D$, Shallit-style); run the
fast-path Cornacchia on each.  5--20 million discriminants per second on 16 cores; yield $\Tt(\sqrt B)$
orders for cost $\Tt(B)$.
\item \emph{Batched smoothness}: maintain each candidate's running $n^4$-smooth part across an adaptive
ladder of cached prime-product segments starting just above $n^2$ (Bernstein remainder trees
\cite{Bern04}; the ladder deepens toward $n^4$ or widens the discriminant pool by a work-balance
trigger, and winners usually appear octaves below $n^4$).
\item \emph{Commit}: for the winner $(D, t)$, compute the class polynomial for the best class invariant
(25--50$\times$ smaller than the Hilbert polynomial \cite{Suth12}) by the CRT method \cite{Suth11},
parallelized across CRT primes; find one root $\bmod\ p$; map invariant${}\to{}j$; assemble the Montgomery
curve and a point of exact order $m$.  Group-structure obstructions (a Montgomery model needs
$4 \mid N$, and additionally $8 \mid N$ when $p \equiv 3 \bmod 4$; a point of order $m$ needs
$m \mid \text{exponent}$) are \emph{removed rather than filtered}: for each prime $\ell$ with non-cyclic
$\ell$-torsion, walk the $\ell$-isogeny volcano to its floor \cite{Suth13,FM02}, where the
$\ell$-Sylow subgroup is cyclic and the curve is always representable.
\end{enumerate}
Because the scan must reach $B \approx (\#\text{candidates})^2$, the end-to-end heuristic cost is
$p^{1/4+o(1)}$ even though only $p^{1/8+o(1)}$ candidates are consumed; $H_D$ and root-finding stay at
$\Ot(p^{1/8})$ ($h(D) \approx \sqrt{|D|} \approx p^{1/8}$; the largest committed class polynomial to
date had $h = 35{,}085$ at $10^{100}$).  Direct enumeration of solutions $(t, v) \to D$ is worse
($\Tt(p^{3/8})$ over the $v$-range that matters), so closing the gap to $p^{1/8}$ would require
enumerating usable $(D, t, v)$ in time proportional to their number --- not currently known, and the
natural open problem of this section.

\subsubsection{The measured SEA/CM crossover}\label{sec:race}
A head-to-head race on the same primes \osf, 192-core machines, everything verified by
\texttt{voneshot.py} (Table~\ref{tab:race}).

\begin{table}[htbp]
\centering\scriptsize
\setlength{\tabcolsep}{4.5pt}
\begin{tabular}{@{}rrrrrrrl@{}}
\toprule
bits & $u/2$ & $1/\rho(u/2)$ & CM wall & SEA wall & CM cands (pool) & SEA cands & SEA vs CM (core-s/cand)\\
\midrule
256 & 4.00 & $2.0\times 10^2$ & 4.1 s & 68 s & 6{,}644 & 728 & --- (16-core box)\\
384 & 5.59 & $1.6\times 10^4$ & 2{,}306 s & 615 s & 126{,}282 & 16{,}000 & 7.2 vs 3.5\\
416 & 5.97 & $4.8\times 10^4$ & 1{,}774 s & 2{,}805 s & 152{,}484 & 52{,}000 & 10.1 vs 2.2\\
448 & 6.36 & $1.6\times 10^5$ & 7{,}182 s & (stopped, 20{,}035 s) & 309{,}264 & ${\ge}\,275{,}600$ & 13.7 vs 4.5\\
\bottomrule
\end{tabular}
\medskip
\caption{The SEA-vs-CM one-shot race \osf.}
\label{tab:race}
\end{table}

Individual walls flip on winner-density luck (the density draw has exponential-scale variance), but the
per-candidate costs are unambiguous: SEA pays a flat 7--14 core-seconds while CM's average ---
mostly-linear segment building plus the quadratically growing scan and winner class polynomial --- was
still 2--4.5 core-seconds at pools up to $3\times 10^5$.  Extrapolating both curves puts the plain-SEA
crossover at roughly \textbf{480--512 bits}; SEA-side accelerations that need SEA internals (an
early-abort gate from the order modulo small $\ell$ during the Elkies phase, batch tails, torsion-family
bias; modeled $\times$3--5) would pull it toward ${\approx}\,380$ bits, while the CM route's scan and
winner-classpoly economics cap it near ${\approx}\,500$ bits regardless of patience.  Beyond the
crossover SEA sampling continues as $p^{1/8+o(1)}$ until the density itself walls out around 640 bits
(${\approx}\,3\times 10^8$ candidates ${}\times{} {\approx}\,20$ core-seconds ${}\approx{}$ two hundred
core-years).  The observed winner densities track the Dickman model within a factor of 1--8 across all
measured sizes --- the ``fudge'' collecting the window/multiplicity/torsion structure conditions and
sampling variance.

\subsubsection{Special forms: algebraic shortcuts}\label{sec:shortcuts}
Two independent contributions to the challenge table exploit structure instead of search
\cite{OneShot}:
\begin{itemize}
\item \emph{The supersingular shortcut} (McLain + Claude, NIST P-384 and P-521).  If
$p \equiv 3 \bmod 4$, the curve $y^2 = x^3+x$ is supersingular with order exactly $p+1$ and is already
Montgomery ($A = 0$); a certificate exists iff $\smoothpart{n^4}{p+1} > L$, with no discriminant search
and no class polynomial.  Generalized-Mersenne primes make this likely because $p+1$ factors
algebraically: for P-384, $p+1 = 2^{32}(2^{64}-1)(2^{288}+2^{224}+2^{160}+2^{96}-1)$ has a 197-bit
smooth part against a 193-bit $L$.  For $p = 2^{521}-1$ (P-521), $p+1 = 2^{521}$ is entirely smooth and
the certificate is a power-of-two Pomerance triple with $m = 2^{261}$.  (For $p \equiv 1 \bmod 4$ the
route is impossible in this format: $p+1 \equiv 2 \bmod 4$ while every Montgomery order is divisible by
4.)  This is the elliptic analogue of the classical $p+1$ tests of \cite{BLS75}.
\item \emph{Reverse-CM order generation} (McLain + Claude Fable 5, $10^{80}+129$, 25 CPU-minutes): the
same architecture as \S\ref{sec:cmoneshot} rediscovered independently at small scale ---
\texttt{qfbcornacchia} over $D \le 8\times 10^6$, an $8 \mid N$ intake filter for Montgomery
representability, budgeted Brent--rho smooth parts, a subset-product search for $m$ with exponents
capped by $v_q(m) \le v_q(N) - \min(\lfloor v_q(N)/2\rfloor, v_q(p-1))$ so that $m$ divides the group
exponent for every possible group structure, and one \texttt{polclass} for the winner.
\end{itemize}

\subsection{Results to date}\label{sec:oneshotresults}

The challenge table \cite{OneShot} holds verified certificates for the least prime above $10^c$ for
$c = 27, \dots, 120$ and for cryptographically famous primes.  Selected entries, all verified by
\texttt{voneshot.py}, appear in Table~\ref{tab:oneshotresults}.

\begin{table}[htbp]
\centering\footnotesize
\begin{tabular}{@{}llll@{}}
\toprule
$p$ & bits & route & time / resources\\
\midrule
$10^{27}+103 \,\dots\, 10^{50}+151$ & 90--167 & SEA sampling (\texttt{oneshot.gp}) & 2--170 CPU-s each\\
$10^{55}+21 \,\dots\, 10^{75}+129$ & 183--250 & CM engine \osf & ${\approx}$1--7 s on 16 cores\\
$2^{255}-19$ (Curve25519) & 255 & CM engine & ${\approx}$4 s on 16 cores\\
P-256; secp256k1 field primes & 256 & CM engine & ${\approx}$11 s; ${\approx}$8 s (16 cores)\\
$10^{80}+129$ & 266 & reverse-CM (independent) & ${\approx}$25 CPU-min\\
$10^{85}+103 \,\dots\, 10^{105}+3$ & 283--349 & CM engine & 24--500 s on 16 cores\\
$10^{110}+7$, $10^{115}+79$, $10^{120}+79$ & 366--399 & CM engine & 121--226 s on 192--256 cores\\
BLS12-381 field prime & 381 & CM engine & ${\approx}$18 s on 256 cores\\
NIST P-384 field prime & 384 & supersingular shortcut & ${<}$1 s, no search\\
$2^{521}-1$ (P-521) & 521 & supersingular; $m = 2^{261}$ & ${<}$1 s, no search\\
\bottomrule
\end{tabular}
\medskip
\caption{Selected one-shot certificates from the challenge table \cite{OneShot}.}
\label{tab:oneshotresults}
\end{table}

The frontier for \emph{generic} primes stands just below 400 bits ($10^{120}+79$, 2 minutes on 256
cores); the race of \S\ref{sec:race} reached 448 bits in two spot-hours, and the density model prices
generic 512--640-bit certificates at tens to hundreds of core-years --- possible, not yet done.
Special-form primes (smooth $p+1$) escape the wall entirely, as P-521 shows at 521 bits.

\section{Restricted one-shot ECPP: $m$ must be $n^2$-smooth}\label{sec:restricted}

\subsection{The problem, and where it sits in the hierarchy}

Now forbid the listed primes: the certificate is $(p, A, x_0, m)$ with \emph{every} prime factor of $m$
below $n^2$ --- the $k = 0$ stratum of \S\ref{sec:oneshot}, equivalently a zero-level short ECPP
(\S\ref{sec:short}).  The verifier is unchanged (\texttt{voneshot.py} with an empty list): it factors
$m$ completely by trial division below $n^2$ and needs no supplied data beyond $m$ itself.  Existence
for every $p > 3$ again follows from Pomerance triples ($2$ is $n^2$-smooth); the certificate is
${\approx}\,2.5n$ bits.  Verification costs $O((n/2)(1+\log k)\,M(n) + M(n^2))$ bit operations for a
certificate whose $m$ has $k$ distinct primes (\S\ref{sec:shortverif}): Pomerance-order $O(n^2\log n)$
for bounded $k$ --- a Pomerance triple is the case $k = 1$ --- and $O(n^2(\log n)^2)$ in the worst case
$k = \Theta(n/\log n)$.  What changes --- dramatically --- is the search.

\subsection{Finding: the exponent doubles}\label{sec:restrictedfind}

The usable-curve condition becomes $\smoothpart{n^2}{\#E} > L$: the same half of $\#E$'s bits must now
be built from primes below $n^2$ instead of $n^4$.  The Dickman parameter doubles
($u = \ln p/\ln n^2 = 2\ln p/(4\ln\ln p)(1+o(1))$), so the per-candidate probability is
\[
  \rho(u/2)^{1+o(1)} \;=\; p^{-1/4+o(1)},
\]
the square of the $n^4$ density: \textbf{$p^{1/4+o(1)}$ candidates via SEA sampling, and
$p^{1/2+o(1)}$ end-to-end for the CM route} (scan budget $B \approx \text{candidates}^2$) --- the CM
engine loses its practical edge entirely here, degenerating to the same $\sqrt p$ order as the
Pomerance brute force, while SEA sampling stays genuinely sub-$\sqrt p$.  Concretely, with the race's
measured SEA rates (Table~\ref{tab:frontier}):

\begin{table}[htbp]
\centering\footnotesize
\begin{tabular}{@{}lrrl@{}}
\toprule
bits & $u/2$ ($y{=}n^2$) & $1/\rho(u/2)$ & SEA cost estimate\\
\midrule
67 ($10^{20}$) & 2.74 & 12 & seconds, one core\\
100 ($10^{30}$) & 3.75 & $1.1\times 10^2$ & seconds, one core\\
133 ($10^{40}$) & 4.71 & $1.3\times 10^3$ & minutes, one core\\
167 ($10^{50}$) & 5.62 & $1.7\times 10^4$ & ${\approx}$1 h, one core\\
200 ($10^{60}$) & 6.52 & $2.5\times 10^5$ & ${\approx}$100 core-h\\
256 & 8.00 & $3.1\times 10^7$ & ${\approx}$5--10 core-years (fleet-scale)\\
320 & 9.61 & ${\approx}\,9\times 10^9$ & ${\approx}\,10^3$--$10^4$ core-years\\
384 & 11.18 & ${\approx}\,3\times 10^{12}$ & out of reach\\
\bottomrule
\end{tabular}
\medskip
\caption{The feasibility frontier of the restricted variant.}
\label{tab:frontier}
\end{table}

The $n^4$ window with its self-certifying listed primes is thus worth a full halving of the exponent
($1/4 \to 1/8$) at zero asymptotic cost to the verifier --- the cleanest illustration in the family of
how a small amount of certificate data (the $q_i$, a few hundred bits) buys an enormous amount of
findability.

\subsection{Results: first generic certificates}\label{sec:restrictedresults}

The only previously existing restricted certificate at challenge scale was the special-form P-521 entry
($m = 2^{261}$, from $p+1 = 2^{521}$).  For this write-up we produced generic restricted certificates
for the least primes above $10^{20}, \dots, 10^{50}$ by SEA sampling (a 15-line wrapper over the
repository's \texttt{oneshot.gp} with the smoothness bound set to $n^2$; first-winner-wins parallel
workers; every certificate verified by \texttt{voneshot.py} with an empty prime list).  Observed
candidate counts sit within an order of magnitude of $1/\rho(u/2)$, as expected given the $O(1)$
structure corrections and geometric variance (Table~\ref{tab:restricted}).

\begin{table}[htbp]
\centering\footnotesize
\begin{tabular}{@{}lrrrll@{}}
\toprule
$p$ & bits & predicted & curves & time / cores & $m$ of the certificate\\
\midrule
$10^{20}+39$ & 67 & 12 & 213 & 33 s, 1 core & $2^3{\cdot}5^3{\cdot}7{\cdot}13{\cdot}37{\cdot}47{\cdot}89$\\
$10^{30}+57$ & 100 & $1.1\times 10^2$ & 117 & 33 s, 1 core & $2{\cdot}7^3{\cdot}43{\cdot}109{\cdot}241{\cdot}251{\cdot}6899$\\
$10^{40}+121$ & 133 & $1.3\times 10^3$ & ${\approx}\,8\times 10^3$ & 17 min, 8 workers & $2{\cdot}17{\cdot}37{\cdot}73{\cdot}197{\cdot}283{\cdot}673{\cdot}2797{\cdot}11827$\\
$10^{50}+151$ & 167 & $1.7\times 10^4$ & ${\approx}\,5.3\times 10^4$ & ${\approx}$2 core-h & $2{\cdot}3{\cdot}7{\cdot}29{\cdot}37{\cdot}211{\cdot}337{\cdot}347{\cdot}449{\cdot}467{\cdot}3527{\cdot}12613$\\
$10^{60}+7$ & 200 & $2.5\times 10^5$ & \multicolumn{3}{l}{not run to completion ($4.8\times 10^3$ curves, no winner); ${\approx}\,10^2$--$10^3$ core-h predicted}\\
\bottomrule
\end{tabular}
\medskip
\caption{First generic restricted one-shot certificates (predicted column is $1/\rho(u/2)$; ``curves''
are aggregates over all workers at the first win).}
\label{tab:restricted}
\end{table}

The four certificates, in the format $(p, A, x_0, m)$ --- each verified True by \texttt{voneshot.py}
and stored in \texttt{certs/restricted/} of \osf:

\medskip
{\footnotesize\ttfamily\raggedright
\noindent $p = 10^{20}{+}39$:\quad A = 19352331390561983207, x0 = 80662701055625636343, m = 14084161000\par
\noindent $p = 10^{30}{+}57$:\quad A = 640307721818372374714786296071, x0 = 821106379164368261893854100371, m = 1341825306264338\par
\noindent $p = 10^{40}{+}121$:\quad A = 9206317171817454415338197777434159186965, x0 = 1798766695016208561785106447263414920568, m = 113982529425281128058\par
\noindent $p = 10^{50}{+}151$:\quad A = 52854763757971328428601223341553948579729822167, x0 = 41413023321438028880723015756658658164929271447635, m = 10372368086928809543163162\par}

\medskip
\noindent The observed-over-predicted ratios ($18, 1.1, 6, 3$) show the same order-of-magnitude scatter as the
\S\ref{sec:race} race densities --- single-draw variance dominates at these pool sizes.  The 256-bit
case (${\approx}\,3\times 10^7$ candidates at ${\approx}$5--10 core-seconds each) is a months-scale
project on a ${\sim}3000$-core fleet with stock SEA --- feasible in principle, unattempted; 384 bits is
out of reach.  For comparison, the \emph{unrestricted} one-shot search handled 256 bits in 4 seconds
(\S\ref{sec:race}): the two exponents $1/8$ vs $1/4$ are separated by exactly this practical chasm.

\section{Short ECPP}\label{sec:short}

\subsection{The problem}\label{sec:shortdef}

The ShortPrimalityProofs repository \cite{Short} adds the one ingredient the one-shot format forbids ---
recursion --- while keeping everything else as small as possible.

\begin{definition}[{\cite{Short}}]\label{def:short}
A \emph{short ECPP} is a sequence of integers
\[
  (p_0,\, A_0,\, x_0,\, m_0p_1,\, A_1,\, x_1,\, m_1p_2,\, \dots,\, A_k,\, x_k,\, m_kp_{k+1})
\]
with $k \ge 0$,
$p_0 \ge 5$, and $n = \lceil\log_2 p_0\rceil$ \emph{fixed for the whole chain}, in which: the $p_i$ are
odd with $n^2 < p_{i+1} < \sqrt{p_i}$ for $i < k$ and $p_{k+1} = 1$; each $m_i$ is $n^2$-smooth with
$L_i < m_ip_{i+1} < r_iL_i$, where $L_i = L(p_i)$ and $r_i$ is the least prime divisor of $m_i$; and
$(x_i, y_i)$ is a point of order $m_ip_{i+1}$ on $B_iy^2 = x^3+A_ix^2+x$ \emph{modulo every prime
divisor of $p_i$}.
\end{definition}

Soundness is downward induction on Lemma~\ref{lem:cert}: $p_{k+1} = 1$ makes the last level's order
fully $n^2$-smooth (a restricted one-shot certificate for $p_k$); each earlier level's order
$m_ip_{i+1} > L_i$ proves $p_i$ prime \emph{given} that $p_{i+1}$ is prime.  Note $m_i \ge 2$ always
(else the order would be $p_{i+1} < \sqrt{p_i} < L_i$), so $r_i$ is well defined; and the certificate
stores only the \emph{product} $m_ip_{i+1}$ --- the verifier recovers $p_{i+1}$ as its $n^2$-rough part
by trial division below $n^2$, and never needs a primality test on it (the next level is the primality
proof).  The ``modulo every prime divisor'' phrasing is essential and is exactly the point of the CRT
counterexample of Remark~\ref{rem:crt}, which was found in the course of pinning this definition down.

\emph{Structure of the chain.}  $p_{i+1} < \sqrt{p_i}$ at least halves the bit length per level, and
$p_{i+1} > n^2$ forces termination once $p_i \le n^4$: so $k = O(\log n) = O(\log\log p_0)$ levels, and
the total size is $\sum_i (2+\tfrac12)\operatorname{bits}(p_i) \approx 5n$ bits --- \emph{linear} in
$\log p_0$, like a Pomerance triple, against the quadratic size of a classical chain.  Existence for
every prime $p_0 \ge 5$ is again inherited from Pomerance triples (take $k = 0$,
$m_0 = 2^{\kappa}$); the repository verifies it exhaustively below $2^8$ (201{,}072 short ECPPs; every
prime $5 \le p \le 2^8$ admits one, while $p = 2, 3$ admit none, which is why the format takes
$p_0 \ge 5$).

\subsection{Verification: Pomerance-order bit complexity}\label{sec:shortverif}

The reference verifier \texttt{vsmallECPP.py} processes the levels in order, working modulo $p_i$ at
level $i$ (write $b_i = \lceil\log_2 p_i\rceil$, so $b_{i+1} \le \lceil b_i/2\rceil$ and
$\sum_i b_i < 2n + O(\log n)$; one sieve to $n^2$ serves the whole chain since $n$ is fixed at the top):
it recovers $m_i$ and $p_{i+1} = o_i/m_i$ by trial division of the carried product $o_i$ over the primes
${}\le n^2$, checks the window, minimality, and descent $p_{i+1}^2 < p_i$, checks that
$[o_i](x_i{:}1)$ is a genuine point at infinity, and establishes exactness of the order by the
\S\ref{sec:oneshotverif} divide-and-conquer with $p_{i+1}$ participating as one of the ``primes'' (its
primality being the next level's job).  For the accounting it pays to organize the exactness checks so
that \emph{the big prime is peeled first}: compute $U = [m_i]P$ and check $\gcd(Z_U, p_i) = 1$ --- this
\emph{is} the $(o_i/p_{i+1})$-cofactor condition --- then check $[p_{i+1}]U = O$, then run the
divide-and-conquer only over the primes of $m_i$, applied to $V = [p_{i+1}]P$.  The tree depth then
multiplies only the filler's bits, never the big prime's.  (The reference implementation keeps
$p_{i+1}$ inside the tree; for bounded $k_i$ this changes nothing, and in general it costs at most one
extra factor $\log k_i$ on the $b_i$ term.)

\begin{proposition}[verification cost, bit operations]\label{prop:shortverif}
Let a valid short ECPP for an $n$-bit $p_0$ have levels $i = 0, \dots, k$ with fillers $m_i$ having
$k_i$ distinct prime factors.  With the primes below $n^2$ and their product tree precomputed --- a
certificate-independent table depending only on $n$, built once in $O(n^2(\log n)^2)$ bit operations
and stored in $O(n^2)$ bits --- the certificate is verified in
\[
  O\Bigl(\, \sum_i \bigl[\, b_i + \log m_i\cdot(1 + \log k_i)\,\bigr]\, M(b_i) \;+\; M(n^2) \Bigr)
\]
bit operations: per level, $O(b_i)$ ladder bits for $[o_i]P$ and the peeled big-prime checks plus
$O(\log m_i)$ ladder bits on each of the $\lceil\log_2 k_i\rceil$ levels of the exactness tree, each
ladder bit costing $O(1)$ multiplications $\bmod\ p_i$; globally, one batched trial division of all the
$o_i$ against the prime table at $M(n^2) = O(n^2\log n)$ \cite{Bern04} (the sieve itself is
$O(n^2\log\log n)$).  Consequently, whenever the fillers satisfy
$\log m_i\cdot\log k_i = O(b_i)$ --- in particular whenever $m_i \le n^{O(1)}$, which holds for every
certificate in the challenge table (the largest filler there has 50 bits, and the largest exponent
$\log m_i/\log n$ is $6.11$; see Remark~\ref{rem:tightshort}) and can always be arranged
(2-power fillers, \S\ref{sec:pomexist}) --- \emph{verification costs $O(n^2\log n)$ bit operations: the
same order as verifying a Pomerance triple, exceeded only by a constant factor}.
\end{proposition}

The constant is roughly an order of magnitude with the reference formulas: the two-register ladder
spends 12 multiplications per scalar bit against the doubling-only chain's 5, the level sum
$\sum b_i$ approaches $2n$ against Pomerance's single $(n/2)$-bit chain, and the trial-division pass
adds ${\approx}\,M(n^2) \approx 2n^2\log_2 n$.  In the format's worst case --- a filler of
$\Theta(b_i)$ bits with $\Theta(b_i/\log b_i)$ distinct primes, which the definition permits and which
is exactly what the $k = 0$ stratum ($=$ the restricted one-shot, \S\ref{sec:restricted}) can look like
--- the tree factor is genuine and verification degrades to $\Theta(n^2(\log n)^2)$, the one-shot
bound, as the inclusion of \S\ref{sec:restricted} forces.  Both statements sharpen the repository's
$O((\log p_0)^{2+o(1)})$; the chain structure --- one big prime per level, geometrically shrinking
moduli, small smooth fillers --- is precisely what deletes the one-shot's extra log factor, the same
mechanism as in the fixed-ratio chains of \S\ref{sec:bdescent}.

\begin{remark}[tightened specifications with worst-case $O(n^2\log n)$ verification]
\label{rem:tightshort}
The two potential sources of a second $\log$ factor in Proposition~\ref{prop:shortverif} are distinct
in kind.  (a)~\emph{The exactness tree}: its depth $\log k_i$ multiplies the filler bits, and the
definition permits a filler of $\Theta(b_i)$ bits with $\Theta(b_i/\log b_i)$ distinct primes, for
which the top level costs $\Theta(n^2(\log n)^2)$ --- a worst case over certificates, vacuous for the
ones the finder produces.  (b)~\emph{The trial division}: the per-certificate cost is
$M(n^2) = \Theta(n^2\log n)$ --- Pomerance-order --- but only given the precomputed product tree of
the primes below $n^2$, whose construction costs $\Theta(n^2(\log n)^2)$; a self-contained accounting
carries the extra log here instead.

\emph{Capping the fillers.}  Requiring only $m_i \le n^c$ for a fixed $c$, with the certificate
payload unchanged, removes both sources --- though one must be careful where the trial division
happens.  Direct per-prime trial division of $o_i$ cannot run in $o(n_i^3)$ at a level of $n_i$ bits:
each of the $\pi(n_i^2) \approx n_i^2/(2\ln n_i)$ candidate primes must scan the $(n_i/2)$-bit
$o_i$, $\Omega(n_i)$ apiece.  Instead the verifier splits the recovery in two stages.  First, extract
the smooth parts of \emph{all} levels in one batched pass: reduce the primorial
$P = \prod_{q \le n^2} q$ modulo $\prod_i o_i$ at cost $M(1.44\,n^2) = O(n^2\log n)$, remainder-tree
down to the individual $o_i$ ($\Ot(n)$), and apply the repeated-squaring gcd trick at each level
($\Ot(n_i)$).  Second, factor each extracted $m_i \le n_i^c$ --- now a $c\log_2 n_i$-bit integer ---
directly: trial division to $n_i$ for its small primes, then deterministic Pollard--Strassen
splitting \cite{Pol74,Str77} of the $O(1)$ remaining factors in $(n_i, n_i^2]$, at
$\Ot(n_i^{c/4}+n_i)$ per level.  Everything after the single batched reduction is quasi-linear per
level and sums geometrically, so verification is worst-case $O(n^2\log n)$ --- with the one caveat
that the primorial, a single $1.44\,n^2$-bit integer depending only on $n$ (about 80\,KB at 665
bits), is a per-bit-size precomputation whose cold construction still costs $\Theta(n^2(\log n)^2)$.
A terminal level should then be permitted only once $\sqrt{p_k} \le n_k^{\,c}$, pushing the chain one
polylogarithmically cheap level deeper.  The exactness checks collapse as in
Proposition~\ref{prop:shortverif}: one half-length ladder $V = [p_{i+1}]P$ per level plus scalars
${\le}\,n_i^c$, for $\sum_i(b_i/2 + O(\log^2 n))$ ladder bits and a constant roughly $5\times$
Pomerance's.

\emph{Empirical compliance.}  Against the twenty published chains (70 levels): keeping the
smoothness bound at the top-level $n^2$, as in the current definition, a cap with $c = 7$ retains
every certificate --- the extremal fillers are 46 bits at $10^{50}{+}151$ (level 0,
$m = n^{6.11}$) and 50 bits at $10^{90}{+}289$ (level 0, $n^{5.97}$).  Re-indexing the bounds per
level ($n_i = \lceil\log_2 p_i\rceil$, fillers $n_i^2$-smooth and ${\le}\,n_i^c$) is the cleaner
specification asymptotically but is not backward compatible: 23 of the 70 levels, hitting 17 of the
20 chains, contain a filler prime exceeding $n_i^2$ (deep and terminal levels were built under the
top-level bound), and the per-level cap would need $c \ge 6.5$.  Either way the finder is essentially
unaffected: the $L(1/3)$ channel's fillers are smooth parts of typical size $n_i^{2+o(1)}$, and a cap
only truncates a thin tail of lucky large-filler winners --- a constant-factor effect, not an
asymptotic one.

\emph{Factored fillers.}  If one is willing to change the payload instead, carrying each filler in
factored form removes even the primorial: the verifier checks each listed prime
$q_j \le n^2$ by trial division to $\sqrt{q_j} \le n$, and the $n^2$-roughness of $p_{i+1}$ need
never be checked at all --- a composite $p_{i+1}$ fails at the next level, by the same induction that
already carries the soundness proof.  Verification is then worst-case $O(n^2\log n)$ with $O(n)$
memory and \emph{no} precomputation of any kind, at the price of certificates that are no longer
unique for a given $(p_i, A_i, x_i, o_i)$.

\emph{The uncapped format, all-inclusive.}  Conversely, if one keeps the current definition
unchanged, the worst case is $O(n^2(\log n)^2)$ \emph{with the per-$n$ precomputation included}:
the primorial product tree costs $\Theta(n^2(\log n)^2)$ to build, one reduction against
$\prod_i o_i$ costs $M(n^2) = O(n^2\log n)$, the shared active-branch descent delivering every
level's prime divisors costs $O(n^2\log n)$ more, and the elliptic worst case (a top-level filler of
$\Theta(n)$ bits with $\Theta(n/\log n)$ primes) is $\Theta(n^2(\log n)^2)$ exactly --- as it must
be, since the $k = 0$ stratum is the restricted one-shot.  The one proviso is organizational: the
trial division must be batched \emph{across the whole chain} (build the tree once, reduce it once
against $\prod_i o_i$, share the descent); repeating the $n^2$-scale pass level by level costs an
extra logarithmic factor.  (The reference verifier currently does repeat the pass per level, and
additionally builds its batch products by sequential multiplication, which is quadratic in the batch
size --- $\Tt(n^3)$ asymptotically, though with so small a constant that it is invisible below
$n \sim 10^6$; both are local changes to the implementation, not to the format.  Patched verifiers
realizing the batched organization, with unchanged accept/reject behavior, accompany
\cite{OSF} in \texttt{verifier-batching-pr/}.)  The three regimes, then: the current definition verifies in
$O(n^2(\log n)^2)$ fully self-contained; capped fillers give $O(n^2\log n)$ plus a once-per-$n$
$O(n^2(\log n)^2)$ table; factored fillers give $O(n^2\log n)$ outright.
\end{remark}

One formula in fact covers the elliptic work of the whole family of
\S\S\ref{sec:pomerance}--\ref{sec:short}:
$O(\sum_i [\log o_i + \log m_i(1+\log k_i)]\,M(b_i))$ bit operations, where $o_i$ is the certified
point order at level $i$ and $m_i$ its smooth part with $k_i$ distinct primes.  A Pomerance triple is
one level with $\log m = n/2$ and $k = 1$: $O(n\,M(n)) = O(n^2\log n)$.  A restricted one-shot has
$\log m = n/2$ with $k$ free: Pomerance-order for bounded $k$, and $\Theta(n^2(\log n)^2)$ at
$k = \Theta(n/\log n)$.  The $n^4$ one-shot is the same with its listed primes included in $k$.  A
short ECPP keeps every $\log m_i$ small and pushes all the large bits into the once-per-level big
prime --- which is why it alone combines subexponential findability with Pomerance-order verification.

\subsection{Finding certificates: the free prime makes it subexponential}\label{sec:shortfind}

The search problem at one level of the chain: given $p$ (a level modulus), find a curve and a divisor
$o = mq$ of $\#E$ with $m$ $n^2$-smooth, $q$ prime in $(n^2, \sqrt p)$, and $o$ in the window
$(L, rL)$.  Compare \S\ref{sec:restricted}: the order may now contain \emph{one arbitrary prime} up to
$\sqrt p$ --- but the prover must \emph{discover} $q$, i.e.\ partially factor $\#E$.  The reference
finder \texttt{short.gp} \cite{Short} (SEA on random curves and their twists, so two orders per point
count) recognizes winners through three channels of increasing cost:
\begin{enumerate}
\item \emph{Terminal:} some divisor of the $n^2$-smooth part $s$ of $\#E$ lands in the window --- a
restricted one-shot certificate; ends the chain ($p_{\text{next}} = 1$).  Density as in
\S\ref{sec:restricted}: $\rho(u/2)$ with $u = \ln p/\ln n^2$.
\item \emph{Cheap prime rough part:} the rough cofactor $R = \#E/s$ is itself prime and below
$\sqrt p$.  Costs one pseudoprimality test, but demands that the entire discarded half of $\#E$ be
$n^2$-smooth: measured density 25\% of curves at $10^{10}$, 0 of 200 by $10^{40}$ --- this channel
alone does not scale.
\item \emph{General descent (the workhorse):} factor $R$ under a time budget (\texttt{factorint}
inside an \texttt{alarm}; the budget is an early abort that keeps the search on curves whose orders
factor easily) and use any exposed prime $q \in (n^2, \sqrt p)$ with a smooth filler $m \mid s$
completing the window.  The discarded cofactor $\#E/o$ is unconstrained --- this is what channel (2)
lacked.
\end{enumerate}
Unlucky sub-chains are abandoned by bounded backtracking (a descent gets a budget of 64 curves before
the parent resamples), and \texttt{parallel\_short.py} races independent workers mixing complete- and
partial-factorization strategies.  At small levels complete factorization is available and the finder
descends faster than square root (observed chains like $665 \to 333 \to 159 \to 47$ bits).

\begin{proposition}[heuristic finding cost]\label{prop:shortfind}
Model channel (3) with ECM as the budgeted factoring engine and reach $y$: a candidate order
$N \approx p$ wins if $N = S\cdot q$ with $q$ prime ${}\le \sqrt p \approx N^{\sigma}$
($\sigma = \tfrac12$) and the cofactor $S$ $y$-smooth; winners concentrate at cofactor size
${\approx}\,p^{\sigma}$, so the per-candidate probability is $\rho(u)$,
$u = \sigma\ln p/\ln y$, and each candidate costs an SEA count plus ECM peeling
$L_y(1/2, \sqrt 2)$.  Optimizing $y$ (which lands at $y = \Lp{2/3}{\cdot}$) gives total heuristic cost
\[
  \Lp{1/3}{(9\sigma/2)^{1/3}} \;=\; \Lp{1/3}{(9/4)^{1/3} \approx 1.31}
  \qquad\text{(SEA supply)},
\]
and with a CM-supplied candidate pool (scan cost ${\approx}\,\mathrm{candidates}^2$) the constant
becomes
\[
  (2/3)(18\sigma)^{1/3} \approx 1.39 \qquad\text{(CM supply)}.
\]
Levels shrink geometrically, so the top level dominates and the whole chain costs the same ---
subexponential, with a smaller constant than the $\Lp{1/3}{1.923}$ of factoring $p$ itself by NFS
\cite{LL93}.
\end{proposition}

The derivation and the $b$-parametrized generalization are in \osf{} (report \emph{descent-chain});
the robust conclusions are the exponent $1/3$, the $(1-b)^{1/3}$ shape, and ``cheaper than
factoring.''

\subsection{Results to date}\label{sec:shortresults}

The challenge table \cite{Short} holds verified short ECPPs for the least prime above $10^c$,
$c = 10, 20, \dots, 200$ --- the top entry at 665 bits, twice the one-shot frontier, found on hardware
where the deepest one-shot runs needed hundreds of cores (Table~\ref{tab:shortresults}).

\begin{table}[htbp]
\centering\footnotesize
\begin{tabular}{@{}lrrl@{\qquad}lrrl@{}}
\toprule
$p$ & bits & levels & CPU time & $p$ & bits & levels & CPU time\\
\midrule
$10^{10}+19$  & 34  & 2 & ${<}$1 s        & $10^{110}+7$    & 366 & 4 & ${\approx}$850 s\\
$10^{20}+39$  & 67  & 2 & ${\approx}$1 s  & $10^{120}+79$   & 399 & 4 & ${\approx}$1400 s\\
$10^{30}+57$  & 100 & 2 & ${\approx}$1 s  & $10^{130}+1113$ & 432 & 3 & ${\approx}$55 s\\
$10^{40}+121$ & 133 & 3 & ${\approx}$2 s  & $10^{140}+13$   & 466 & 4 & ${\approx}$800 s\\
$10^{50}+151$ & 167 & 3 & ${\approx}$3 s  & $10^{150}+67$   & 499 & 5 & ${\approx}$1500 s\\
$10^{60}+7$   & 200 & 3 & ${\approx}$3 s  & $10^{160}+303$  & 532 & 5 & ${\approx}$1200 s\\
$10^{70}+33$  & 233 & 4 & ${\approx}$45 s & $10^{170}+63$   & 565 & 4 & ${\approx}$830 s\\
$10^{80}+129$ & 266 & 4 & ${\approx}$180 s& $10^{180}+313$  & 599 & 4 & ${\approx}$850 s\\
$10^{90}+289$ & 299 & 3 & ${\approx}$350 s& $10^{190}+253$  & 632 & 4 & ${\approx}$1400 s\\
$10^{100}+267$& 333 & 3 & ${\approx}$800 s& $10^{200}+357$  & 665 & 4 & ${\approx}$4400 s\\
\bottomrule
\end{tabular}
\medskip
\caption{The ShortPrimalityProofs challenge table \cite{Short}.  Entries through $c = 100$ by
\texttt{short.gp} on a single core (AVS and Claude Code, Fable 5); $c = 110, \dots, 200$ by
\texttt{parallel\_short.py} workers.}
\label{tab:shortresults}
\end{table}

The growth from 3 s at 200 bits to 4{,}400 s at 665 bits --- a factor ${\approx}\,10^3$ over a
$\times 3.3$ size increase, with large per-draw variance (the lucky 432-bit chain cost 55 s) --- is
consistent with the gentle $L(1/3)$ curve; extrapolating, $10^{300}$ (997 bits) is a
core-days-to-weeks project.

\subsection{Interpolating toward classical ECPP: fixed-ratio descent chains}\label{sec:bdescent}

Relaxing the descent rate from ``square root'' to a parameter: a chain in which each step's certified
prime satisfies $q \le p^b$ (and the point order is just $c\cdot q$ with no minimality demanded) was
analyzed and prototyped in \osf{} (report \emph{descent-chain}) for $b \in (\tfrac12, 1)$: certificate
size $O(n)/(1-b)$; verification $O(n^2\log n)/(1-b^2)$ bit operations --- Pomerance's order, the
geometric shrink killing the one-shot's extra log factor; finding
$\Lp{1/3}{(2/3)(18(1-b))^{1/3}}$, which tends to $0$ as $b \to 1$ --- recovering polynomial-time
classical ECPP in the limit, at the price of $\Theta(n)$ levels and quadratic certificates.  Short ECPP
is the $b = \tfrac12$ endpoint with minimized orders.  Measured with $b = 0.85$ on a 16-core desktop:
$2^{255}-19$ proved in 3.7 s (vs ${\approx}$15 s one-shot), $10^{100}+267$ in 3.7 s (vs 7.9 min
one-shot), 1024 bits in 128 s (9 steps, verification 405 ms), 2048 bits in 5.2 h (12 steps,
verification 2.65 s) --- sizes far beyond any one-shot's reach, with verification still two orders of
magnitude cheaper than a classical chain's at equal size.

\section{Classical ECPP}\label{sec:ecpp}

\subsection{The problem and the two classical algorithms}\label{sec:ecppalgs}

The classical elliptic curve primality proof drops all structural demands on the order: a chain
$p = p_0 > p_1 > \cdots$, where one step exhibits, for $p = p_i$, a curve $E/\F_p$ with
$\#E = c\cdot q$, $q$ a probable prime exceeding $L(p)$, together with a point $P$ such that
$Q = [c]P \ne O$ and $[q]Q = O$; the next level certifies $q = p_{i+1}$.  By Lemma~\ref{lem:cert}
(applied to the point $Q$ of prime order $q > L(p)$ modulo every prime divisor) each step is sound, and
the chain bottoms out at a size testable by trial division.  The exhibited order is a single prime; the
cofactor $c$ needs no factorization and carries no burden of proof --- its peelability is purely the
prover's recognition problem.  Two ways to run a step:

\begin{itemize}
\item \textbf{SEA supply (Goldwasser--Kilian \cite{GK86}; the ``basic'' route):} sample random curves
over $\F_p$, count points, and keep a curve with $\#E = 2q$, $q$ a probable prime (${\approx}\,p/2$, so
the chain sheds one bit per step).  Expected $O(\log p)$ curves per step by the prime number theorem
\emph{if} the Hasse interval contains the expected number of primes --- provable for all inputs outside
a set of primes of relative density zero using unconditional results on primes in short intervals
\cite{HB88}, conjecturally for all.  With SEA at $\Ot((\log p)^4)$ per count and $\Theta(\log p)$
curves per step, a step costs $(\log p)^{5+\eps}$, and the downrun --- $\Theta(\log p)$ steps for the
strict $m = 2q$ descent, a logarithmic factor fewer with polynomially smooth cofactors peeled ---
puts the route at $(\log p)^{6+o(1)}$ in total.
Adleman and Huang \cite{AH92} closed the gap of exceptional primes by a random reduction through
Jacobians of genus-2 curves, putting primality in ZPP: \emph{every} prime has an
elliptic-curve-style certificate findable in expected polynomial time --- a theorem, though never a
practical algorithm.
\item \textbf{CM supply (Atkin--Morain \cite{AM93}; fastECPP):} replace point counting by the CM supply
of \S\ref{sec:cmengine}: for each small discriminant $D$ with $4p = t^2+|D|v^2$, the candidate orders
$p+1\pm t$ are known before any curve is built; trial-divide them, keep one of the form $c\cdot q$ with
$c$ smooth and $q$ a probable prime, and only then construct the curve from a root of $H_D$.  Because
tiny discriminants suffice (the step needs only \emph{some} usable order, an event of probability
${\approx}\,1/\log p$ per candidate), the class polynomials stay small, and the certified prime can be
pushed down by the cofactor's length (typically $\Theta(\log n)$ bits per step, giving
$\Theta(n/\log n)$ to $\Theta(n)$ steps).  Heuristic cost $(\log p)^{5+\eps}$ for the basic version ---
dominated by the per-discriminant square roots --- and $(\log p)^{4+\eps}$ for \emph{fastECPP}
\cite{FKMW04,Mor07}, which batches the square roots $\sqrt D \bmod p$ across all $D$ via one square
root per small prime discriminant factor (Shallit's trick): the fastest known heuristic primality
\emph{proving} complexity for general $p$.  Enge's MPI implementation \cite{Enge24} parallelizes it to
quasi-cubic wall-clock time on a linear number of cores.
\end{itemize}

\subsection{Certificates and verification}\label{sec:ecppverif}

A classical certificate stores, per step,
$(p_i, \text{curve}, m_i, q_{i+1}, P_i)$: ${\approx}\,5$ integers of $b_i$ bits with $b_i$ shrinking
only $\theta = \Theta(\log n)$ per step, so $\Theta(n^2/\log n)$ bits in all.  Verification re-runs the
two scalar multiplications of every step: $\sum_j O(b_j)\,M(b_j)$ over $\Theta(n/\theta)$ steps of
linearly decaying size is $O(n^3\log n/\theta) = O(n^3)$ bit operations for the customary cofactor
lengths --- the $M$-factor cancels against the step count --- quasi-cubic, the benchmark against which
the $O(n^2\log n)$-class formats of \S\S\ref{sec:pomerance}--\ref{sec:short} are the improvement.
(Verification is embarrassingly parallel across steps, a fact used in practice for the record
certificates below.)

\subsection{Results to date}\label{sec:ecppresults}

Classical ECPP is the workhorse of general-purpose primality proving, with three decades of records:
from the sub-1{,}000-digit proofs of the first ECPP implementations \cite{AM93} through
${\approx}\,10^4$--$2\times 10^4$ digits in the fastECPP era of the mid-2000s \cite{FKMW04,Mor07}, to
today's frontier held by Enge's CM/fastECPP-over-MPI \cite{Enge24} and Marcel Martin's Primo (the
standard desktop tool).  The largest ECPP-proven primes as of mid-2026 \cite{t5k}:

\begin{center}\footnotesize
\begin{tabular}{@{}lrll@{}}
\toprule
prime & digits & date & notes\\
\midrule
$R(109297) = (10^{109297}-1)/9$ & 109{,}297 & May 2025 & repunit; A.~Enge \& P.~Underwood, fastECPP/MPI\\
$R(86453)$ & 86{,}453 & May 2023 & repunit; previous record\\
$5^{104824}+1048245$ & 73{,}269 & Feb 2023 & \\
$(2^{221509}-1)/292391881$ & 66{,}673 & Oct 2023 & Mersenne cofactor\\
\bottomrule
\end{tabular}
\end{center}

For perspective across the whole write-up: at $10^5$ digits the \emph{only} practical route is the
classical one --- every other format in this survey buys its cheaper verification with a finding cost
that is subexponential-to-exponential in the bit length and hence confined to hundreds of bits
(one-shot) or ${\approx}\,10^3$ bits (short/descent chains).  Conversely, at the sizes where they do
apply, the \S\ref{sec:pomerance}--\ref{sec:short} formats verify orders of magnitude faster than a
classical chain, with certificates orders of magnitude smaller.

\section{Synthesis and open problems}\label{sec:synthesis}

The five formats populate a two-parameter design space: \emph{how much order structure is demanded} (a
power of two; $n^2$-smooth; $n^4$-smooth with listed witnesses; smooth${}\times{}$one-prime; a single
free prime) and \emph{how much recursion is allowed} (none; square-root descent; ratio-$b$ descent;
free descent).  Moving along either axis trades finding cost against verification cost and certificate
size, and Table~\ref{tab:summary} summarizes the endpoints.  The empirical record across all four
challenge repositories matches the heuristic model remarkably well: candidate counts within a factor
1--20 of the Dickman/counting predictions at every measured size, with the systematic part of the
discrepancy accounted for by known local biases and structure conditions (\S\ref{sec:supply}), and
single-draw geometric variance dominating the rest.

Open problems, in roughly increasing order of depth:
\begin{enumerate}
\item \emph{Find a Pomerance triple for $10^{27}+103$} \cite{DANGER3} --- the standing challenge;
${\approx}\,3\times 10^{13}$ candidates at the measured GPU rate is around ten GPU-hours.
\item \emph{Beat $\Tt(\sqrt p)$ for Pomerance triples}, or give evidence none of the standard toolkits
(CM, isogenies, torsion families, sieving) can.  The obstruction is that ``$2^k$ divides $\#E$'' is a
codimension-$(n/2)$ condition with no multiplicative structure (\S\ref{sec:pomfind}).
\item \emph{Close the one-shot $p^{1/8}$ vs $p^{1/4}$ gap}: enumerate usable $(D, t, v)$ with
$4p = t^2+|D|v^2$ in time proportional to their number rather than to the discriminant bound $B$.  Any
such enumeration would make the CM route $p^{1/8+o(1)}$ end-to-end and move the SEA crossover
(\S\ref{sec:race}) beyond reach.
\item \emph{Instrument SEA for early-abort supply}: expose the order modulo small $\ell$ during the
Elkies phase to gate candidates before the full $\Ot(n^4)$ count (modeled $\times$3--5 at one-shot
sizes; needs internals no current library exports).
\item \emph{A generic 256-bit restricted certificate} (\S\ref{sec:restrictedresults}):
${\approx}\,3\times 10^7$ SEA candidates --- months on a ${\sim}3000$-core fleet with stock tools, much
less with item 4.
\item \emph{Rigorous existence with the right densities.}  Existence of all four quasi-quadratic
formats is unconditional (via \cite{Pom87}), but every \emph{finding} bound in this write-up is
heuristic; even the classical chain's polynomial time is only ``almost all $p$'' provable
\cite{GK86,HB88} (all $p$ in expected polynomial time only via \cite{AH92}).  Proving that a positive
proportion of curves over every $\F_p$ has $\smoothpart{y}{\#E} > \sqrt p$ for the relevant $y$ would
derandomize a whole column of the table.
\item \emph{Lower bounds.}  Pomerance's discussion \cite[\S1]{Pom87} still delimits what is known: some
primes (Fermat, Mersenne) have $(1+o(1))\log_2 p$-multiplication proofs; nothing rules out
$O(\log\log p)$ multiplications for special primes; and whether $(5/2)\log_2 p$ is optimal for general
$p$ --- he declined to conjecture it --- remains untouched.  The formats here show the \emph{finding}
cost of the known constructions rising steeply as verification approaches his constant; whether that
trade-off is intrinsic is open.
\item \emph{The three-way trade-off.}  Exhibit a certificate family achieving simultaneously:
finding in (even randomized, even heuristic) polynomial time; size $n^{1+o(1)}$; and verification in
$O(n^{2+o(1)})$ bit operations.  Short ECPP achieves the last two with $\Lp{1/3}{{\approx}1.31}$
finding; the AKS--Bernstein certificate (Appendix~\ref{app:aks}) achieves the first two --- finding in
expected \emph{quadratic} time --- with $(\log p)^{4+o(1)}$ verification, a quartic floor that is
structural for its proof technique.  No known family achieves all three.
\end{enumerate}

\appendix

\section{The AKS--Bernstein certificate}\label{app:aks}

This appendix spells out, from \cite{Ber07}, the certificate summarized in Remark~\ref{rem:aks}: the
definition, the existence theorem with witness set $S = \{1\}$, the verification algorithm with its
cost accounting, and the comparison with the elliptic formats.  Throughout, $p > 1$ is the integer to
be certified and $n = \lceil\log_2 p\rceil$.

\subsection{The certificate}

\begin{definition}[{\cite[Def.~3.1]{Ber07}}]\label{def:akscert}
Let $d, e$ be positive integers, $c, c_-$ integers, $f$ a monic polynomial in $(\Z/p\Z)[y]$ of degree
$d$, $R = (\Z/p\Z)[y]/f$, $r \in R$, and $S \subseteq R$ a finite set.  Then
$(d, e, c, c_-, f, r, S)$ is a \emph{certificate} for $p$ if:
\begin{enumerate}
\item $e$ divides $p^d - 1$, and $e > c \ge c_- \ge 0$;
\item $r^{p^d-1} = 1$ in $R$, and $r^{(p^d-1)/q} - 1$ is a unit in $R$ for each prime $q \mid e$;
\item every $s \in S$ is a unit in $R$; $s^e - (s')^e$ is a unit for all distinct $s, s' \in S$; and
$s^e - r$ is a unit for all $s \in S$;
\item $\displaystyle\binom{e\#S}{c_-}\binom{c}{c_-}\binom{e\#S - c_- + e-1-c}{e-1-c} \;\ge\;
p^{d\sqrt{e/3}}$; and
\item $(x-s)^{p^d} = r^{(p^d-1)/e}\,x - s$ in the ring $R[x]/(x^e - r)$, for every $s \in S$.
\end{enumerate}
If $p$ has a certificate then $p$ is a power of a prime \cite[Thm.~3.2]{Ber07}.
\end{definition}

Nothing is assumed about $f$ beyond monicity --- in particular the verifier never tests
irreducibility; soundness maps $R$ onto a field over a prime divisor of $p$ regardless.  The proof of
soundness is an AKS-type counting argument: condition (2) makes $r^{(p^d-1)/e}$ an exact $e$-th root
of unity $\zeta$ locally, condition (5) lets Frobenius-like maps generate at least
$p^{d\sqrt{e/3}}$ distinct products of the $\zeta^i x - s$ (this is where condition (4) enters), and a
Minkowski-lattice collision among the maps $u \mapsto u^{p^i P^j}$ then forces $p$ to be a prime
power.  A $(\log p)^{1+o(1)}$ perfect-power test \cite{Ber98} upgrades ``prime power'' to
``prime.''

\subsection{Existence: $S = \{1\}$ always suffices}

\begin{theorem}[{\cite[Thms.~5.1--5.3]{Ber07}}]\label{thm:akscert}
Let $p$ be prime.
\begin{enumerate}
\item For every $d \ge 1$, every divisor $e \ge 6$ of $p^d-1$ with $e \ge d^2\lceil\log_2 p\rceil^2$,
every monic \emph{irreducible} $f$ of degree $d$, and every $r \in R$ such that $r^{(p^d-1)/e}$ has
order $e$: the tuple $(d, e, 0, 0, f, r, \{1\})$ is a certificate for $p$.
\item Such a pair $(d, e)$ exists for every $p \ge 2$, and the least admissible $d$ satisfies
$d \le \exp(O(\log\log\log p \cdot \log\log\log\log p)) = (\log p)^{o(1)}$, so that
$e = (\log p)^{2+o(1)}$; this rests on lower bounds for the smooth part of $p^d - 1$ of
Odlyzko--Pomerance type.
\end{enumerate}
Consequently every prime has a certificate with $\#S = 1$ and $c = c_- = 0$, of total size
$(\log p)^{1+o(1)}$ bits: the small integers $d, e$, and the elements $f, r$ of
$O(d\log p)$ bits each.
\end{theorem}

Irreducibility of $f$ and the order condition on $r$ enter only here, to \emph{guarantee} existence
($R$ is then a field, and any generator of $R^*$ works for $r$); an $f$ that happens to be reducible
can still appear in a valid certificate.  The finder of \cite[\S4]{Ber07} samples both: an irreducible
$f$ in expected $d$ trials, and a valid $r$ in expected $\prod_{q \mid e} q/(q-1) = (\log p)^{o(1)}$
trials, each trial one exponentiation of cost $(\log p)^{2+o(1)}$; locating $(d, e)$ costs
$(\log p)^{2+o(1)}$ by batched remainder trees.  Total expected finding time:
$(\log p)^{2+o(1)}$.  Larger witness sets $S$ permit smaller $e$ --- the gain in condition (4) is only
logarithmic in $\#S$, however, while the cost of condition (5) grows linearly in $\#S$, so bounded
$\#S$ is asymptotically optimal and the freedom is used for constant-factor tuning
\cite[\S7]{Ber07}.

\subsection{The verification algorithm}\label{sec:aksverif}

Given $p$ and a claimed certificate $(d, e, c, c_-, f, r, S)$ with $d, \#S \le (\log p)^{o(1)}$ and
$e \le (\log p)^{2+o(1)}$ (``reasonably small'' in \cite[\S6]{Ber07}; the verifier rejects larger
parameters):
\begin{enumerate}
\item Check that $p$ is not a perfect power \cite{Ber98}.  \hfill $(\log p)^{1+o(1)}$
\item Compute $p^d - 1$; check $e \mid p^d - 1$ and $e > c \ge c_- \ge 0$.  \hfill
$(\log p)^{1+o(1)}$
\item Check condition (4): the three binomial coefficients are integers of $(\log p)^{2+o(1)}$ bits,
computable by balanced products; compare against $p^{d\sqrt{e/3}}$ via $(\log p)^{2+o(1)}$-bit
arithmetic.  \hfill $(\log p)^{2+o(1)}$
\item Check $r^{p^d-1} = 1$ in $R$: $d\log_2 p$ squarings in $R$, each $(\log p)^{1+o(1)}$ bit
operations.  \hfill $(\log p)^{2+o(1)}$
\item Factor $e$ by trial division; for each of its $(\log p)^{o(1)}$ primes $q$, check that
$r^{(p^d-1)/q} - 1$ is a unit in $R$ (a unit test in $R$ is a polynomial gcd with $f$ plus integer
gcds with $p$, $(\log p)^{1+o(1)}$).  \hfill $(\log p)^{2+o(1)}$
\item Check the unit conditions on $S$: each $s$, each $s^e - r$, each $s^e - (s')^e$.  \hfill
$(\log p)^{1+o(1)}$ per test
\item \emph{The big exponentiation.}  For each $s \in S$, compute $(x-s)^{p^d}$ in
$R[x]/(x^e - r)$ by square-and-multiply and compare with $r^{(p^d-1)/e}x - s$.  A ring element
occupies $e \cdot d \cdot \log_2 p = (\log p)^{3+o(1)}$ bits, one ring multiplication costs
$(\log p)^{3+o(1)}$ (Kronecker substitution), and there are $d\log_2 p = (\log p)^{1+o(1)}$
squarings: \hfill $(\log p)^{4+o(1)}$
\end{enumerate}
Accept iff all checks pass; then $p$ is a prime power by Definition~\ref{def:akscert}, hence prime by
step 1.  Every step is deterministic and unconditional.  The total is $(\log p)^{4+o(1)}$ bit
operations, entirely dominated by step 7.

\subsection{Why quartic, and the comparison with the elliptic formats}\label{sec:aksquartic}

The quartic cost of step 7 is structural, not an artifact: condition (4) --- the engine of the
soundness proof --- forces $e\#S \gtrsim \sqrt{e}\,d\log p$ up to logarithmic factors, i.e.\
$e = \widetilde\Omega((\log p)^2)$ for bounded $\#S$, so that the congruence ring has
$\widetilde\Omega((\log p)^3)$ bits; and the exponent $p^d$ forces $\Omega(\log p)$ ring squarings.
Within this proof technique, $(\log p)^{4+o(1)}$ verification is therefore a floor (larger $\#S$
trades a polylogarithmically smaller $e$ against linearly more congruences, and loses).  The upshot
for this survey's central trade-off: the AKS--Bernstein certificate achieves \emph{polynomial ---
indeed expected quadratic --- finding} and \emph{linear size}, but quartic verification; the short
ECPP of \S\ref{sec:short} achieves \emph{linear size} and \emph{quasi-quadratic verification}, but
$\Lp{1/3}{\cdot}$ finding.  Neither dominates the other, and no known certificate family combines
polynomial finding, linear size, and quasi-quadratic verification --- the closing open problem of
\S\ref{sec:synthesis}.

\begin{thebibliography}{FKMW04}

\bibitem[AH92]{AH92} L.~M. Adleman, M.-D.~A. Huang, \emph{Primality Testing and Abelian Varieties over
Finite Fields}, Lecture Notes in Math.\ 1512, Springer, 1992.

\bibitem[AKS04]{AKS04} M.~Agrawal, N.~Kayal, N.~Saxena, PRIMES is in P, \emph{Ann.\ of Math.}\ 160
(2004), 781--793.

\bibitem[AM93]{AM93} A.~O.~L. Atkin, F.~Morain, Elliptic curves and primality proving, \emph{Math.\
Comp.}\ 61 (1993), 29--68.

\bibitem[Bac90]{Bach90} E.~Bach, Explicit bounds for primality testing and related problems,
\emph{Math.\ Comp.}\ 55 (1990), 355--380.

\bibitem[Ber04]{Bern04} D.~J. Bernstein, How to find smooth parts of integers, preprint, 2004,
\url{https://cr.yp.to/papers.html#smoothparts}.

\bibitem[Ber07]{Ber07} D.~J. Bernstein, Proving primality in essentially quartic random time,
\emph{Math.\ Comp.}\ 76 (2007), 389--403.

\bibitem[Ber98]{Ber98} D.~J. Bernstein, Detecting perfect powers in essentially linear time,
\emph{Math.\ Comp.}\ 67 (1998), 1253--1283.

\bibitem[Berr05]{Berr05} P.~Berrizbeitia, Sharpening ``PRIMES is in P'' for a large family of
numbers, \emph{Math.\ Comp.}\ 74 (2005), 2043--2059.

\bibitem[BLS75]{BLS75} J.~Brillhart, D.~H. Lehmer, J.~L. Selfridge, New primality criteria and
factorizations of $2^m\pm 1$, \emph{Math.\ Comp.}\ 29 (1975), 620--647.

\bibitem[BLS12]{BLS12} R.~Br\"oker, K.~Lauter, A.~V. Sutherland, Modular polynomials via isogeny
volcanoes, \emph{Math.\ Comp.}\ 81 (2012), 1201--1231.

\bibitem[CEP83]{CEP83} E.~R. Canfield, P.~Erd\H{o}s, C.~Pomerance, On a problem of Oppenheim concerning
``factorisatio numerorum,'' \emph{J.~Number Theory} 17 (1983), 1--28.

\bibitem[Cor08]{Cor08} G.~Cornacchia, Su di un metodo per la risoluzione in numeri interi
dell'equazione $\sum C_h x^{n-h}y^h = P$, \emph{Giornale di Matematiche di Battaglini} 46 (1908),
33--90.

\bibitem[Cox13]{Cox13} D.~A. Cox, \emph{Primes of the Form $x^2+ny^2$}, 2nd ed., Wiley, 2013.

\bibitem[CP05]{CP05} R.~Crandall, C.~Pomerance, \emph{Prime Numbers: A Computational Perspective}, 2nd
ed., Springer, 2005.

\bibitem[Deu41]{Deu41} M.~Deuring, Die Typen der Multiplikatorenringe elliptischer
Funktionenk\"orper, \emph{Abh.\ Math.\ Sem.\ Univ.\ Hamburg} 14 (1941), 197--272.

\bibitem[Dic30]{Dic30} K.~Dickman, On the frequency of numbers containing prime factors of a certain
relative magnitude, \emph{Ark.\ Mat.\ Astr.\ Fys.}\ 22A (1930), 1--14.

\bibitem[Elk98]{Elk98} N.~D. Elkies, Elliptic and modular curves over finite fields and related
computational issues, in \emph{Computational Perspectives on Number Theory}, AMS/IP Stud.\ Adv.\ Math.\
7 (1998), 21--76.

\bibitem[Eng09]{Enge09} A.~Enge, The complexity of class polynomial computation via floating point
approximations, \emph{Math.\ Comp.}\ 78 (2009), 1089--1107.

\bibitem[Eng24]{Enge24} A.~Enge, FastECPP over MPI, 2024, \href{https://arxiv.org/abs/2404.05506}
{arXiv:2404.05506}; software CM/cm-ecpp, \url{https://www.multiprecision.org/cm/ecpp.html}.

\bibitem[FKMW04]{FKMW04} J.~Franke, T.~Kleinjung, F.~Morain, T.~Wirth, Proving the primality of very
large numbers with fastECPP, in \emph{Algorithmic Number Theory (ANTS-VI)}, LNCS 3076, Springer, 2004.

\bibitem[FM02]{FM02} M.~Fouquet, F.~Morain, Isogeny volcanoes and the SEA algorithm, in
\emph{Algorithmic Number Theory (ANTS-V)}, LNCS 2369, Springer, 2002, 276--291.

\bibitem[GK86]{GK86} S.~Goldwasser, J.~Kilian, Almost all primes can be quickly certified, in
\emph{Proc.\ 18th STOC}, 1986, 316--329; journal version: Primality testing using elliptic curves,
\emph{J.~ACM} 46 (1999), 450--472.

\bibitem[Gra08]{Gran08} A.~Granville, Smooth numbers: computational number theory and beyond, in
\emph{Algorithmic Number Theory}, MSRI Publ.\ 44, Cambridge, 2008, 267--323.

\bibitem[HB88]{HB88} D.~R. Heath-Brown, The number of primes in a short interval, \emph{J.~reine
angew.\ Math.}\ 389 (1988), 22--63.

\bibitem[Hil86]{Hil86} A.~Hildebrand, On the number of positive integers $\le x$ and free of prime
factors $> y$, \emph{J.~Number Theory} 22 (1986), 289--307.

\bibitem[HvdH21]{HvdH21} D.~Harvey, J.~van der Hoeven, Integer multiplication in time $O(n\log n)$,
\emph{Ann.\ of Math.}\ 193 (2021), 563--617.

\bibitem[Koh96]{Kohel96} D.~Kohel, \emph{Endomorphism rings of elliptic curves over finite fields},
PhD thesis, University of California, Berkeley, 1996.

\bibitem[Len87]{Len87} H.~W. Lenstra, Jr., Factoring integers with elliptic curves, \emph{Ann.\ of
Math.}\ 126 (1987), 649--673.

\bibitem[LL93]{LL93} A.~K. Lenstra, H.~W. Lenstra, Jr.\ (eds.), \emph{The Development of the Number
Field Sieve}, Lecture Notes in Math.\ 1554, Springer, 1993.

\bibitem[LP19]{LP19} H.~W. Lenstra, Jr., C.~Pomerance, Primality testing with Gaussian periods,
\emph{J.~Eur.\ Math.\ Soc.}\ 21 (2019), 1229--1269.

\bibitem[McL26]{McL26} A.~McLain (with GPT 5.5 Codex / Claude), Danger 2026 Data Challenge code and
research notes, 2026, \url{https://github.com/alexamclain/Danger2026DataChallenge}.

\bibitem[Mil76]{Mil76} G.~L. Miller, Riemann's hypothesis and tests for primality, \emph{J.~Comput.\
System Sci.}\ 13 (1976), 300--317.

\bibitem[Mon87]{Mon87} P.~L. Montgomery, Speeding the Pollard and elliptic curve methods of
factorization, \emph{Math.\ Comp.}\ 48 (1987), 243--264.

\bibitem[Mor07]{Mor07} F.~Morain, Implementing the asymptotically fast version of the elliptic curve
primality proving algorithm, \emph{Math.\ Comp.}\ 76 (2007), 493--505.

\bibitem[Pol74]{Pol74} J.~M. Pollard, Theorems on factorization and primality testing,
\emph{Proc.\ Cambridge Philos.\ Soc.}\ 76 (1974), 521--528.

\bibitem[Pom87]{Pom87} C.~Pomerance, Very short primality proofs, \emph{Math.\ Comp.}\ 48 (1987),
315--322.

\bibitem[Pra75]{Pra75} V.~R. Pratt, Every prime has a succinct certificate, \emph{SIAM J.~Comput.}\ 4
(1975), 214--220.

\bibitem[Rab80]{Rab80} M.~O. Rabin, Probabilistic algorithm for testing primality, \emph{J.~Number
Theory} 12 (1980), 128--138.

\bibitem[Rue26]{Rue26} F.~Ruehle (with Claude Opus 4.6), Danger 2026 Data Challenge implementation,
2026, \url{https://github.com/ruehlef/Danger2026DataChallenge}; see also J.~Shi,
\url{https://github.com/JaneShi99/pomerance-p24}.

\bibitem[Sch85]{Scho85} R.~Schoof, Elliptic curves over finite fields and the computation of square
roots mod $p$, \emph{Math.\ Comp.}\ 44 (1985), 483--494.

\bibitem[Sch87]{Scho87} R.~Schoof, Nonsingular plane cubic curves over finite fields, \emph{J.~Combin.\
Theory Ser.~A} 46 (1987), 183--211.

\bibitem[Sch95]{Scho95} R.~Schoof, Counting points on elliptic curves over finite fields,
\emph{J.~Th\'eor.\ Nombres Bordeaux} 7 (1995), 219--254.

\bibitem[Sil09]{Sil09} J.~H. Silverman, \emph{The Arithmetic of Elliptic Curves}, 2nd ed., GTM 106,
Springer, 2009.

\bibitem[SS14]{SS14} I.~E. Shparlinski, A.~V. Sutherland, On the distribution of Atkin and Elkies
primes, \emph{Found.\ Comput.\ Math.}\ 14 (2014), 285--297.

\bibitem[Str77]{Str77} V.~Strassen, Einige Resultate \"uber Berechnungskomplexit\"at,
\emph{Jahresber.\ Deutsch.\ Math.-Verein.}\ 78 (1976/77), 1--8.

\bibitem[Sut11]{Suth11} A.~V. Sutherland, Computing Hilbert class polynomials with the Chinese
remainder theorem, \emph{Math.\ Comp.}\ 80 (2011), 501--538.

\bibitem[Sut12]{Suth12} A.~V. Sutherland, Accelerating the CM method, \emph{LMS J.~Comput.\ Math.}\ 15
(2012), 172--204.

\bibitem[Sut13]{Suth13} A.~V. Sutherland, Isogeny volcanoes, in \emph{ANTS-X}, Open Book Ser.\ 1, MSP,
2013, 507--530.

\bibitem[t5k]{t5k} The PrimePages, Top twenty: Elliptic Curve Primality Proof,
\url{https://t5k.org/top20/page.php?id=27} (accessed Aug 2026).

\bibitem[Wat69]{Wat69} W.~C. Waterhouse, Abelian varieties over finite fields, \emph{Ann.\ Sci.\
\'Ecole Norm.\ Sup.}\ (4) 2 (1969), 521--560.

\bibitem[DANGER3]{DANGER3} A.~V. Sutherland, \emph{DANGER3} (Pomerance triples: definition, verifier
\texttt{vpp.py}, Lean verifier, exhaustive data to $2^{32}$, challenge log),
\url{https://github.com/AndrewVSutherland/DANGER3}, 2026.

\bibitem[OneShot]{OneShot} A.~V. Sutherland, \emph{OneShotPrimalityProofs} (one-shot ECPP: definition,
verifier \texttt{voneshot.py}, finder \texttt{oneshot.gp}, exhaustive data to $2^8$, challenge table),
\url{https://github.com/AndrewVSutherland/OneShotPrimalityProofs}, 2026.  Part of the DARPA expMath
program.

\bibitem[Short]{Short} A.~V. Sutherland, \emph{ShortPrimalityProofs} (short ECPP: definition, verifier
\texttt{vsmallECPP.py}, finders \texttt{short.gp}/\texttt{parallel\_short.py}, exhaustive data to
$2^8$, challenge table), \url{https://github.com/AndrewVSutherland/ShortPrimalityProofs}, 2026.

\bibitem[OSF]{OSF} A.~V. Sutherland and Claude Code, \emph{OneShotFastECPP} (the CM one-shot
engine: \texttt{dscan}/\texttt{smooth}/\texttt{fproot}/\texttt{cm\_method}/\texttt{oneshot}; reports
\emph{descent-chain}, \emph{sea-crossover}, and this write-up; verified certificates),
\url{https://github.com/AndrewVSutherland2/OneShotFastECPP}, 2026.

\end{thebibliography}

\end{document}
